\documentclass[11pt]{article}
\PassOptionsToPackage{obeyspaces,spaces}{url}
\usepackage[margin=1in]{geometry}
\usepackage{amsmath,amssymb}
\usepackage{graphicx}
\usepackage{booktabs}
\usepackage{longtable}
\usepackage{url}
\usepackage[colorlinks=true,linkcolor=blue,citecolor=blue]{hyperref}
\usepackage{caption}
\captionsetup{font=small}

\graphicspath{{figs/}{docs/figs/}{../model_output/moist_v2_sweep/figs/}%
{/home/shill/py/axisym-hc-swm/model_output/moist_v2_sweep/figs/}}

\newcommand{\Hhat}{\hat{H}}
\newcommand{\Shat}{\hat{S}}
\newcommand{\Wstar}{W^{*}}
\newcommand{\Wq}{W_q}
\newcommand{\py}{\partial_y}
\newcommand{\dyr}{\Delta_\theta^{\rm rad}}
\newcommand{\mmday}{mm\,day$^{-1}$}
\newcommand{\ms}{m\,s$^{-1}$}

\title{A parameter sweep of the moist V2 shallow-water Hadley model\\
\large 149 simulations across the sign change of the gross moist stability,
with off-equatorial and seasonally migrating forcing}
\author{Run and written by Claude, for Spencer Hill}
\date{2026-08-09}

\begin{document}
\maketitle

\tableofcontents

\section{Summary of what the runs show}
\label{sec:summary}

151 runs, of which 150 completed: 131 with time-invariant forcing and 19 with a
seasonally migrating forcing, all at $n_y = 801$ except the resolution twins, all
integrated 5800 days (5.01 drag timescales) except the seasonal ones and one
short pair. One configuration diverged and was replaced.

\begin{enumerate}

\item \textbf{One number organises most of it.} $r = \Wq/\Wstar$, the ratio of
the moisture an undisturbed column holds to the moisture at which the gross
moist stability changes sign, is the control parameter
(Section~\ref{sec:control}). Four runs with $a$ from 0.75 to 0.87 and $W_c$ from
34.9 to 52.0, chosen so that all four sit at $r = 0.85$, agree on the
subtropical jet to 0.19\%, on peak $|v|$ to 0.10\%, on the raining fraction to
0.00\% and on peak precipitation to 1.02\%. Across the 58 runs that vary the
gross moist stability by any route, those same quantities span 77\%, 382\%, 59\%
and 839\% of their means. The collapse is real, not a restatement of the
definition.

\item \textbf{Two of the six parameters do not collapse, for identifiable
reasons.} $\tau_c$ and $E_0$ each move $r$ and also move something else: how
fast convection removes water, and how much water there is. The $\tau_c$ ladder
crosses $r = 1$ (0.906 to 1.084) with the domain raining everywhere at every
rung, while a run at $r = 1.055$ reached through $W_c$ rains over only 73\% of
the domain.

\item \textbf{The transition begins before $\Hhat$ changes sign.} On the finest
transect ($a = 0.85$, nine values of $W_c$) the domain rains everywhere up to
$r = 0.919$ and has split into three bands over 91\% of the domain by
$r = 0.964$. Convergence raises $W$ above $\Wq$ in the ascending region, so the
local $\Hhat$ goes negative there before the undisturbed column does.

\item \textbf{What the aggregated state does with its energy.} At $a = 0.85$,
$W_c = 50$ the mean circulation carries $+104.7$~MW\,m$^{-1}$ into the
ascending column, because $\Hhat < 0$ there, and the eddy moisture flux carries
$-106.2$ back out. They cancel to 1.4\%. The eddy flux is the brake, and the
solution finds its equilibrium by drawing part of the domain dry enough to turn
$\Hhat$ positive: over $1.0 < |y| < 3.8$~Mm the column sits at
43.7~kg\,m$^{-2}$ against $\Wstar = 44.25$, giving $\Hhat = +1.0$~MJ\,m$^{-2}$
where the rest of the domain has $-33$ (Section~\ref{sec:budget}).

\item \textbf{Removing the eddy moisture flux removes the brake, completely.}
At $D = 0$ on the negative-stability reference, 4.6\% of the domain rains, at
831.7~\mmday, in 35 separate bands. At $D = 4\times10^{6}$~m$^2$\,s$^{-1}$,
71.8\% rains at a peak of 34.4. On the positive-stability reference the same
range of $D$ changes peak precipitation by 3\%.

\item \textbf{Aggregation needs the feedback at nearly full strength.} With
$\Lambda$ at 0, 0.25, 0.50 and 0.75 of its physical value $L_v/C$ the
negative-stability reference rains everywhere, peaking at 4.9, 5.2, 5.7 and
7.9~\mmday. At the physical value it rains over 63\% of the domain, peaking at
77.8. The onset between 0.75 and 1.0 is abrupt.

\item \textbf{With passive moisture the same parameter does nothing at all.}
Two V1 twins, one at $r = 0.76$ and one at $r = 1.14$, have identical
circulations: jet 28.20~\ms, peak $|v|$ 0.365~\ms, raining everywhere, and no
organised band. Their gross moist stabilities are $+18.4$ and
$-11.4$~MJ\,m$^{-2}$. $\Hhat$ is a diagnostic in V1 and a dynamical quantity in
V2, and this is the control that shows it.

\item \textbf{There is a second boundary, further out, where the solutions stop
settling.} 18 of the 131 never reach a steady state. In 16 the jet magnitude
itself swings, by 2.6 to 83.0~\ms{} over the trailing 2000 days, along with the
peak rain rate by up to 2890~\mmday. In the other two the jet holds to under
0.7~\ms{} while the rain bands wander across the domain, which shows up as a
pointwise $u$ difference of 10\% and 62\% between consecutive 1000-day means.
They span $r = 1.132$ to 2.290, and the largest $r$ that still settles is 1.310,
so unsteadiness sets in well past $\Hhat = 0$ rather than at it
(Section~\ref{sec:numerics}).

\item \textbf{The initial moisture does not select the attractor.} Starting $W$
15 or 30~kg\,m$^{-2}$ above $W_c$ at four points spanning the transition changes
the equilibrium by at most $8\times10^{-4}$ relative. No multiple equilibria
were found (Section~\ref{sec:multistable}).

\item \textbf{The rain band and the energy flux equator are 1.9 Mm apart.} With
the forcing peak at 1400~km, the rain peaks at 1.082~Mm and the total column MSE
flux vanishes at 2.959~Mm. Because $\Hhat$ varies only from 17.8 to
19.3~MJ\,m$^{-2}$ across the tropics there, the flux is essentially $v\Hhat$ and
vanishes where $v$ does, which is the boundary between the cross-equatorial cell
and the summer cell rather than the ascent maximum
(Section~\ref{sec:offeq}).

\item \textbf{The seasonal response is quasi-static, and a single-timescale
picture of it is wrong by a factor of ten.} A 37-day thermal relaxation predicts
the rain band should lag the forcing by 17 to 37~days depending on period, and
be damped to 0.36 of the quasi-steady amplitude at a 90-day period. Measured
across five periods from 90 to 1440 days, the lag is 3.1 to 6.0~days and the
gain 0.85 to 0.94. The gain also matches the steady off-equatorial value at the
same displacement to within 2\%. These cycles repeat exactly: consecutive cycles
differ by 0.0~km in band position (Section~\ref{sec:seasonal}).

\item \textbf{Every seasonal run on the negative-stability reference is
chaotic.} Consecutive cycles differ by 10,889 to 14,656~km in rain-band
position, so their amplitudes and lags describe one arbitrary realisation and
are not reported as results.

\item \textbf{None of it looks much like the observed Hadley circulation, and
the jet is the worst of it.} Across the plane the subtropical jet runs 1.53 to
2.56 times the observed 31.0~\ms, and with the forcing 2800~km off the equator
the winter jet reaches 128.7~\ms{} with $-54$~\ms{} of equatorial easterly. The
eddy energy transport is essentially absent, as the ERA5 calibration already
found (Section~\ref{sec:plausible}).

\item \textbf{The numerics are sound where it matters and not everywhere.}
Halving $\Delta t$ changes $u$ by 0.01\% of its peak. Across all eight
resolution and timestep twins the jet moves by at most 0.30\% and peak
precipitation by at most 0.91\%. But the profile at $|y| \approx 7$ to 11~Mm is
not converged: between $n_y = 801$ and 401 it differs by up to 67.6\% of peak
$u$, and between 801 and 1601 by up to 9.8\%. The tropics are converged; the
westerly terminus region is not.

\end{enumerate}

\paragraph{Independent check.} Two of these configurations were run in July 2026
by a different session with different analysis code, and their equilibria were
recorded in \path|CLAUDE.md| and in \path|docs/era5_calibration_report.pdf|. All
18 quantities recorded there are reproduced, most to under 1\%: peak
precipitation 77.82 against 78, the secondary band at 5.592~Mm and 18.72~\mmday{}
against 5.6 and 18.7, the dry fraction 0.362 against 0.36, the terminus notch at
7.482~Mm and $-45.49$~\ms{} against 7.48 and $-45$, and the jet 73.79 against
73.8. The quiescent column sits 71.22~K above its radiative target against the
analytic $\tau\Lambda E_0 = 71.22$~K. Run
\path|python scripts/v2_sweep_verify.py| to regenerate that comparison.

\section{The gross moist stability, and the one number that sets its sign}
\label{sec:control}

\subsection{What the quantity is}

The moist V2 model carries column moist static energy
$\langle h\rangle = C\theta + L_v W$, where $C$ is the column heat capacity that
converts a potential-temperature tendency into a column-energy tendency and
$L_v$ is the latent heat of vaporisation. Adding the thermodynamic equation
(\ref{eq:theta}) times $C$ to the moisture equation (\ref{eq:w}) times $L_v$
gives the budget (\ref{eq:mse}), whose transport term is
\begin{equation}
\py\Big[\,\underbrace{v\big(\Shat - L_v(2a-1)W\big)}_{\text{mean circulation}}
\;\underbrace{-\;L_v D\,\py W}_{\text{eddies}}\Big].
\end{equation}
The bracket multiplying $v$ is the gross moist stability
$\Hhat = \Shat - L_v(2a-1)W$. It is the column energy the mean circulation
carries per unit of upper-level meridional velocity, and it is the difference of
two terms that partly cancel. The first, the gross dry stability
$\Shat = C\,d\,\Delta_z/H$, is positive: the upper branch is warmer than the
lower branch, so poleward flow aloft exports dry static energy. The second is
negative: the lower branch is moister than the upper branch and flows the other
way, so the same circulation imports latent energy.

Its \emph{sign} is what matters, and here is why. Suppose ascent at some
latitude strengthens. Ascent means upper-level divergence, $\py v > 0$. If
$\Hhat > 0$, the strengthened circulation exports more energy from that column
than it imports, the column cools, and the ascent is damped. If $\Hhat < 0$ the
strengthened circulation is a net \emph{import} of energy into the ascending
column, which warms and moistens it further, and the ascent grows. A negative
gross moist stability is a positive feedback on ascent. In a full model this is
the mechanism behind self-aggregation of convection; here it is available in a
two-line budget.

\subsection{The two moisture values that decide the sign}

$\Hhat$ vanishes at
\begin{equation}
\Wstar = \frac{\Shat}{L_v(2a-1)} = \frac{C\,d\,\Delta_z}{H\,L_v(2a-1)} .
\label{eq:wstar}
\end{equation}
At the sweep's fixed values ($C = 5.1619\times10^{6}$~J\,m$^{-2}$\,K$^{-1}$,
$d = 4$~km, $\Delta_z = 60$~K, $H = 16$~km, $L_v = 2.5\times10^{6}$~J\,kg$^{-1}$)
this gives $\Shat = 7.7429\times10^{7}$~J\,m$^{-2}$ and
$\Wstar = 44.25$~kg\,m$^{-2}$ at the spec's $a = 0.85$, rising to
$53.40$ at $a = 0.79$ and $61.94$ at $a = 0.75$.

The other value is where the moisture actually sits. In a column with no
transport, the steady moisture balance is $E_0 = P = (W - W_c)/\tau_c$, so
\begin{equation}
\Wq = W_c + \tau_c E_0 .
\label{eq:wq}
\end{equation}
At the defaults $W_c = 50$~kg\,m$^{-2}$, $\tau_c = 14400$~s and
$E_0 = 4.6\times10^{-5}$~kg\,m$^{-2}$\,s$^{-1}$ the correction is
$\tau_c E_0 = 0.66$~kg\,m$^{-2}$, so $\Wq = 50.66$~kg\,m$^{-2}$. This value is
set locally, by evaporation against precipitation, and the circulation cannot
move it except by transporting water in or out faster than convection removes
it. That is why the quiescent parts of the domain sit at $\Wq$ whatever the rest
of the solution is doing.

\subsection{One dimensionless number, five knobs}

Putting (\ref{eq:wstar}) and (\ref{eq:wq}) together, an undisturbed column has
negative gross moist stability when $\Wq > \Wstar$, that is when
\begin{equation}
\boxed{\;
r \;\equiv\; \frac{\Wq}{\Wstar}
\;=\; \frac{L_v(2a-1)\,(W_c + \tau_c E_0)}{C\,d\,\Delta_z/H}
\;>\; 1 . \;}
\label{eq:r}
\end{equation}
Five separate model parameters appear in $r$, and each of them can be moved
across $r = 1$ on its own: the lower-branch water-vapour fraction $a$, the
precipitation threshold $W_c$, the convective timescale $\tau_c$, the
evaporation rate $E_0$, and the static stability $\Delta_z$. That is what makes
the collapse onto $r$ a real question rather than a definition. If $r$ were
merely a convenient relabelling, then runs at the same $r$ reached through
different parameters would still differ. The sweep therefore includes eight runs
in which $W_c$ was chosen for each of four values of $a$ so that all four land
on the same $r$, at $r = 0.85$ and again at $r = 1.05$
(Section~\ref{sec:regime}), and separate ladders in $\tau_c$, $E_0$ and
$\Delta_z$ that cross $r = 1$ by themselves.

Two caveats on $r$, stated now because they limit how much it can be expected
to explain. First, $r$ describes an \emph{undisturbed} column. In the deep
tropics the circulation moves water, so the local $W$ departs from $\Wq$ and the
local $\Hhat$ from its quiescent value; $r$ predicts the sign of the feedback
that gets things started, not the equilibrium reached. Second, $a$ and
$\Delta_z$ appear elsewhere in the equations as well as in $r$. The factor
$(2a-1)$ is also the moisture transport velocity in (\ref{eq:w}), and $\Delta_z$
also multiplies the adiabatic term in (\ref{eq:theta}), so changing either does
more than move $r$. $W_c$, $\tau_c$ and $E_0$ are cleaner: they enter the
dynamics only through $r$ and through the precipitation rate itself.

\subsection{Where the real atmosphere sits}

This is not an exotic corner of parameter space. The ERA5 calibration
(\path|docs/era5_calibration_report.pdf|, its stability section) gives
$\Hhat/\Shat$ between $0.258$ over $|\varphi| \le 10^\circ$ and $0.318$ over
$|\varphi| \le 30^\circ$, so the observed band-mean gross moist stability is
positive but only about a quarter to a third of the dry value: the two competing
terms cancel to within 70\% or so. The same report finds the \emph{local}
$\Hhat$ falling through zero within about $8^\circ$ of the equator. And its
estimate of $\Wstar$ itself, 44 to 57~kg\,m$^{-2}$ depending on how the time
mean is taken, straddles the model's quiescent $50.66$~kg\,m$^{-2}$. So $r = 1$
is roughly where the observations put the tropical atmosphere, and the question
of what happens on either side of it is a question about the real Hadley
circulation.

\section{Where the regime boundary sits}
\label{sec:regime}

\subsection{The map}

Figure~\ref{fig:map} places 25 runs in the plane of $a$ and $W_c$, the two
parameters that move $r$ without touching anything else in the equations. The
line $r = 1$ is drawn on it, and it separates the two behaviours cleanly. Below
it the whole domain rains, at close to the evaporation rate, with a single broad
maximum. Above it the rain concentrates into a few narrow bands separated by
regions where precipitation is identically zero, and the peak rate rises by more
than an order of magnitude.

\begin{figure}[htbp]
\centering
\includegraphics[width=\textwidth]{01_regime_map.png}
\caption{The regime map. Colour shows four properties of the equilibrium; the
solid line is $r = 1$, where an undisturbed column has $\Hhat = 0$, and the
dotted line is $r = 0.85$. Circles settle to a steady state, squares never do.}
\label{fig:map}
\end{figure}

\subsection{The transect through the boundary}

The finest sampling is at $a = 0.85$, where nine values of $W_c$ from 30 to 60
put $r$ between 0.693 and 1.371:

\begin{center}
\begin{tabular}{rrrrrrr}
\toprule
$W_c$ & $r$ & raining fraction & bands & peak $P$ (\mmday) & min $\Hhat$ (MJ\,m$^{-2}$) & jet (\ms)\\
\midrule
30 & 0.693 & 1.000 & 1 & 5.30  & $+23.4$ & 51.5\\
35 & 0.806 & 1.000 & 1 & 6.15  & $+14.4$ & 55.1\\
40 & 0.919 & 1.000 & 1 & 9.02  & $+4.8$  & 60.8\\
42 & 0.964 & 0.910 & 3 & 15.02 & $-0.5$  & 64.7\\
44 & 1.009 & 0.775 & 3 & 38.04 & $-10.7$ & 70.5\\
46 & 1.055 & 0.730 & 5 & 56.64 & $-19.6$ & 74.1\\
50 & 1.145 & 0.633 & 5 & 77.82 & $-32.8$ & 73.8\\
55 & 1.258 & 0.561 & 9 & 97.67 & $-47.3$ & 74.5\\
60 & 1.371 & 0.501 & 9 & 139.52& $-68.0$ & 78.3\\
\bottomrule
\end{tabular}
\end{center}

The boundary is bracketed between $r = 0.919$ and $r = 0.964$, so it sits
\emph{below} $\Hhat = 0$ for the undisturbed column. The reason is in the
caveat of Section~\ref{sec:control}: $r$ describes a column with no transport,
and the ascending region has transport. Moisture convergence lifts $W$ there
above $\Wq$, so the local $\Hhat$ reaches zero while the quiescent collar still
has $\Hhat > 0$. At $W_c = 42$, $r = 0.964$, the minimum $\Hhat$ over the domain
is already $-0.5$~MJ\,m$^{-2}$.

\subsection{The collapse, and its two exceptions}

Figure~\ref{fig:collapse} plots nine properties of the equilibrium against $r$
for every run that varies the gross moist stability by any route. Four routes
collapse: $a$, $W_c$, $\Delta_z$, and the matched-$r$ block that fixes $r$ and
varies $a$ with $W_c$ compensating.

The matched-$r$ block is the test that could have failed. Eight runs, four at
$r = 0.85$ and four at $r = 1.05$, with $a$ from 0.75 to 0.87 and $W_c$ chosen
for each to hit the target $r$:

\begin{center}
\small
\begin{tabular}{lrrrrr}
\toprule
 & raining fraction & peak $P$ & jet & peak $|v|$ & min $\Hhat$\\
\midrule
$r = 0.85$, spread across four $a$ & 0.00\% & 1.02\% & 0.19\% & 0.10\% & 2.71\%\\
$r = 1.05$, spread across four $a$ & 0.34\% & 14.91\% & 2.27\% & 0.07\% & 41.65\%\\
\midrule
for scale: spread over all 58 GMS runs & 58.6\% & 838.6\% & 77.5\% & 382.3\% & 1877.9\%\\
\bottomrule
\end{tabular}
\end{center}

Each spread is the full range divided by the mean. At $r = 0.85$ the four runs
are the same solution to a fraction of a percent, while $a$ differs by 16\% and
$W_c$ by 49\% between them. At $r = 1.05$ the agreement is looser and
structured: the circulation still matches (peak $|v|$ to 0.07\%, jet to 2.3\%)
while the precipitation extremes do not (peak $P$ to 15\%, min $\Hhat$ to 42\%).
Past the boundary the solution has narrow bands whose peak values are sensitive
to where exactly they sit, and $r$ predicts the circulation better than it
predicts the extremes.

\begin{figure}[htbp]
\centering
\includegraphics[width=\textwidth]{02_collapse_on_r.png}
\caption{Nine properties of the equilibrium against $r$, for every run that
varies the gross moist stability. Marker shape and colour identify which
parameter was varied. Filled markers settle to a steady state; open markers
never do, so their plotted value is a 1000-day time mean of a fluctuating state.}
\label{fig:collapse}
\end{figure}

The two exceptions are informative rather than a failure of the idea.
$\tau_c$ enters $r$ through $\Wq = W_c + \tau_c E_0$, and it also sets how fast
convection removes water from a column. Raising it from 1800 to 172800~s at
$a = 0.85$, $W_c = 40$ takes $r$ from 0.906 to 1.084, across the boundary, and
the domain rains everywhere at every rung, peak precipitation only rising from
8.1 to 17.0~\mmday. A slow convective sink smears the rain out, which is the
opposite of what concentrating it requires, so the two effects work against each
other. $E_0$ likewise enters $r$ only weakly, through the same product, while
setting the total water supply outright: doubling it on the negative-stability
reference takes peak $|v|$ from 5.86 to 11.15~\ms{} and peak precipitation from
77.8 to 197.5~\mmday{} at nearly unchanged $r$ (1.145 to 1.160).

$\Delta_z$ collapses for the energetics and the raining fraction, and sits about
6\% low in the jet at matched $r$, because it also multiplies the adiabatic term
in the thermodynamic equation. That was the expected behaviour of the two
parameters flagged in Section~\ref{sec:control} as doing more than moving $r$.

\section{What the solutions on either side look like}
\label{sec:regimes}

\begin{figure}[htbp]
\centering
\includegraphics[width=\textwidth]{03_transect_profiles.png}
\caption{The $a = 0.85$ transect at equilibrium, $W_c$ from 30 to 60, which is
$r$ from 0.693 to 1.371. Labels on the $W$ panel sit in the quiescent collar,
where $W = W_c + \tau_c E_0$ and the curves are furthest apart. The two largest
$W_c$ never settle, and their large extratropical oscillations in $u$ are
discussed in Section~\ref{sec:numerics}.}
\label{fig:transect}
\end{figure}

\subsection{Below the boundary: a Hadley cell that drizzles}

At $r < 0.92$ the solution is a recognisable two-cell Hadley circulation with a
broad rain maximum at the equator sitting on a background that rains everywhere
at the evaporation rate. $W$ is nearly flat, pinned close to $\Wq$, with a bump
of a few kg\,m$^{-2}$ at the equator. Peak precipitation is 5 to 9~\mmday{}
against the 3.97 the uniform evaporation supplies, so the equatorial maximum is
only a 30 to 130\% enhancement over the background. The jet strengthens smoothly
with $r$, from 51.5 to 60.8~\ms{} over the transect, and peak $|v|$ from 1.10 to
2.27~\ms.

\subsection{Above it: bands, dry gaps, and rain rates no atmosphere has}

At $r > 1$ the picture is different in kind. At $W_c = 50$ the equator rains at
77.8~\mmday{} in a spike a few hundred km wide, two secondary bands sit at
$\pm5.59$~Mm raining 18.7~\mmday, and precipitation is identically zero over
36.2\% of the domain. As $r$ rises further the number of bands goes 3, 5, 9 and
the peak rate reaches 139.5~\mmday{} at $r = 1.371$ and 831.7 with the eddy
moisture flux switched off entirely.

Two arithmetic facts make this less mysterious than it first looks, and both are
exact. Peak precipitation is $(W_{\max} - W_c)/\tau_c$, so it is a direct
read-out of the wettest column and nothing more. And because transport moves no
total water, the domain mean of $P$ equals $E_0$ at equilibrium, which the runs
satisfy to between 0.999 and 1.000 in the finite-volume measure. Aggregation
therefore cannot change how much it rains in total; it can only change over how
little of the domain that rain falls. At $W_c = 50$ the raining 63.3\% of the
domain averages 1.58 times $E_0$, and the excess is concentrated in the spike.

The precipitation rates are far outside anything observed. A zonal-mean
77.8~\mmday{} is roughly ten times the wettest zonal-mean band on Earth, and
831.7 is not a rain rate at all. This is the clearest sense in which the
negative-stability solutions are not candidate Hadley circulations.

\section{The solutions against ERA5}
\label{sec:plausible}

Before any of the parameter dependences matter, the question is whether these
are plausible Hadley circulations. The short answer is that the positive
stability side is the right order of magnitude and too strong by roughly a factor
of two, the negative stability side is not a Hadley circulation, and the eddy
energy transport is missing from all of it.

\begin{center}
{\small
\begin{tabular}{cccccccc}
\toprule
$a$ & $W_c$ & $r$ & jet & peak $|v|$ & $\Delta\theta(16^\circ)$ & $(24^\circ)$ & $(33^\circ)$\\
 & & & (m\,s$^{-1}$) & (m\,s$^{-1}$) & (K) & (K) & (K)\\
\midrule
\multicolumn{3}{l}{ERA5} & 31.0 & 1.07--1.31 & 0.56 & 3.21 & 8.75\\
\midrule
0.75 & 30 & 0.49 & 47.5 & 0.83 & 0.38 & 2.34 & 8.31\\
0.79 & 40 & 0.76 & 53.6 & 1.27 & 0.56 & 3.33 & 11.33\\
0.85 & 40 & 0.92 & 60.8 & 2.27 & 0.92 & 5.07 & 15.96\\
0.85 & 44 & 1.01 & 70.5 & 5.03 & 1.60 & 7.33 & 20.91\\
0.85 & 50 & 1.14 & 73.8 & 5.86 & 1.71 & 7.62 & 21.62\\
\multicolumn{3}{l}{V1 twin, $a$=0.79 $W_c$=40} & 28.2 & 0.37 & 0.18 & 1.13 & 4.16\\
\bottomrule
\end{tabular}
}

\end{center}

\medskip
\noindent The ERA5 row is the 2000--2009 annual mean from
\path|docs/era5_calibration_report.pdf| (its Table~9), with peak $|v|$ quoted as
a slab velocity at the two branch depths that report estimates, because that is
the quantity the model's $v$ is. Temperature contrasts are
$\theta(0) - \theta(y)$ at the Coriolis-mapped latitudes $y = 2$, 3 and 4~Mm.

\paragraph{The jet is too strong everywhere, by 1.53 to 2.56 times.} Over the
whole $(a, W_c)$ plane the subtropical jet runs between 47.5 and 79.2~\ms{}
against an observed 31.0. The best case in the sweep is the driest,
$a = 0.75$, $W_c = 30$ at $r = 0.495$, which gives 47.5~\ms. The ERA5
calibration already traced this overshoot to the radiative retarget rather than
to the moisture parameters, and the attribution chain here confirms it
quantitatively: the V1 twin has a 28.20~\ms{} jet, retargeting the relaxation to
$\dyr = 75$~K with the feedback severed raises it to 41.58 with no change in
moisture physics at all, and the sensitivity is $0.708$~\ms{} per K of $\dyr$ on
the positive-stability reference. Getting the jet to 31~\ms{} at fixed
everything-else would take $\dyr \approx 43$~K, which is below the
radiative-convective 50~K and so is not available.

\paragraph{The temperature contrast is too steep by 1.5 to 2.4 times}, in the
same direction and for the same reason. At $24^\circ$ the observed
$\Delta\theta$ is 3.21~K; the runs give 2.34 at $r = 0.49$, 3.33 at 0.76, 5.07
at 0.92 and 7.62 at 1.14. The one run that matches the observed contrast closely
is the V1 twin at 1.13~K, which is too flat.

\paragraph{Column water vapour is about right on the dry side.} ERA5 gives
39.6~kg\,m$^{-2}$ averaged over $|\varphi| \le 20^\circ$. The runs give 40.8 at
$r = 0.76$ and 41.0 at $r = 0.92$, so the reference states hold close to the
observed tropical column. That is unsurprising, since $W$ is pinned near $\Wq$
and $\Wq$ was chosen near observed values, but it does mean the moisture is not
the thing that is wrong.

\paragraph{The energy transport is the largest gap, and it is structural.}
The column fluxes at $24^\circ$, which is $y = 3$~Mm:

\begin{center}
{\small
\begin{tabular}{cccccc}
\toprule
$a$ & $W_c$ & $r$ & mean $v\hat H$ & eddy $-L_vD\py W$ & total\\
 & & & (MW\,m$^{-1}$) & (MW\,m$^{-1}$) & (MW\,m$^{-1}$)\\
\midrule
\multicolumn{3}{l}{ERA5} & 27.9 & 57.7 & 85.6\\
\midrule
0.75 & 30 & 0.49 & 31.54 & 0.16 & 31.70\\
0.79 & 40 & 0.76 & 21.25 & 0.38 & 21.63\\
0.85 & 40 & 0.92 & 11.24 & 0.50 & 11.74\\
0.85 & 44 & 1.01 & 4.86 & -0.65 & 4.21\\
0.85 & 50 & 1.14 & 2.29 & -2.19 & 0.10\\
\multicolumn{3}{l}{V1 twin, $a$=0.79 $W_c$=40} & 6.73 & 0.09 & 6.82\\
\bottomrule
\end{tabular}
}

\end{center}

\medskip
\noindent At $24^\circ$ ERA5 transports 85.6~MW\,m$^{-1}$ of column moist static
energy northward, of which 57.7 is eddy and 27.9 is the mean circulation. The
positive-stability reference gives 21.3 mean and 0.38 eddy. So the mean flux is
within a factor of one and a half of the observed mean flux, and the eddy flux is
smaller than observed by a factor of 150. The reason is the one the ERA5 report
identified: the model's only eddy energy flux is $-L_v D\,\py W$, there is no
eddy dry-static-energy flux in the equations at all, and $W$ is nearly flat
because the local evaporation-precipitation balance pins it, so there is almost
no gradient for $D$ to act on. Raising $D$ to 4 million m$^2$\,s$^{-1}$, above
the top of the ERA5-calibrated range, moves the eddy flux from 0.38 to
1.42~MW\,m$^{-1}$.

The negative-stability side inverts the problem rather than fixing it: at
$a = 0.85$, $W_c = 50$ the eddy flux at $24^\circ$ is $-2.19$~MW\,m$^{-1}$,
pointing the wrong way, because $W$ decreases away from the ITCZ there and the
downgradient flux is equatorward.

\section{The column energy budget}
\label{sec:budget}

\begin{figure}[htbp]
\centering
\includegraphics[width=0.86\textwidth]{04b_anatomy_negative.png}
\caption{How the negative-stability solution at $a = 0.85$, $W_c = 50$
($r = 1.145$) reaches equilibrium. Top to bottom: column water vapour against
the two thresholds that matter, precipitation on a log scale, the gross moist
stability with the non-raining regions shaded, and the column MSE flux resolved
into its mean and eddy parts.}
\label{fig:anatomy}
\end{figure}

\subsection{Never the net alone}

The mean column MSE flux $v\Hhat$ is the difference of two much larger terms,
and reporting it alone would misrepresent the transport by a large factor. At
$24^\circ$ on the positive-stability reference the dry-static flux $\Shat v$ and
the latent flux $-L_v(2a-1)vW$ nearly cancel, and the ERA5 calibration measured
the same cancellation in the observations (its Table~8: $+50.7$ against $-22.9$
observed, $+114.5$ against $-100.5$ in the model at $a = 0.77$). The model
over-transports each component and under-transports their difference.

\subsection{How a negative-stability solution closes its budget}

Figure~\ref{fig:anatomy} answers this for $a = 0.85$, $W_c = 50$, and the answer
has two parts.

\textbf{The eddy flux cancels the mean flux at the ITCZ, to 1.4\%.} At the
ascending column the mean circulation carries $+104.7$~MW\,m$^{-1}$ northward on
one side of the spike and the same southward on the other, importing energy
because $\Hhat < 0$ there. The eddy moisture flux carries $-106.2$, exporting it
down the enormous $W$ gradient the spike creates. The total is
$-1.5$~MW\,m$^{-1}$, 1.4\% of either term. This is what stops the positive
feedback: the ascent intensifies until the moisture gradient it builds is steep
enough for $D$ to remove the energy the mean circulation is importing. That is
also why $D$ is the parameter with the most leverage on this state
(Section~\ref{sec:ladders}).

\textbf{The circulation dries part of the domain back above $\Hhat = 0$.} Over
$1.0 < |y| < 3.8$~Mm the column is drawn down to a mean 43.7~kg\,m$^{-2}$,
below $\Wstar = 44.25$, so $\Hhat$ there is $+1.0$~MJ\,m$^{-2}$ against $-32.8$
at the ITCZ, reaching $+1.70$ at $2.28$~Mm. Those bands are the only part of the
domain from which the mean circulation can export energy, and they lie inside
the non-raining region.

That second part is weaker than the picture suggests at a glance, and the number
is worth stating rather than the impression: 43\% of the non-raining area has
$\Hhat > 0$, and averaged over the whole non-raining area $\Hhat$ is
$-2.54$~MJ\,m$^{-2}$. The two broad gaps between the ITCZ and the secondary
bands are positive; the narrow gaps outside the secondary bands are not. So the
drying is part of the answer and the eddy cancellation is most of it.

\subsection{Budget closure as an equilibrium check}

For the 113 runs that settle, the residual of
$\py(F_{\rm mean} + F_{\rm eddy}) = Q$ is between $3.4\times10^{-9}$ and
$6.3\times10^{-4}$ of the maximum source, and 110 of them are below $10^{-6}$.
The three above it are identifiable: the two members of the off-equatorial
resolution pair, which were run 2500 days rather than 5800 and so are not fully
equilibrated ($5\times10^{-6}$ and $2\times10^{-6}$), and the $\tau_c = 1800$~s
run at $6.3\times10^{-4}$, whose $u$ drift is $2\times10^{-5}$ so it is settled;
its residual is presumably the diagnostic reconstruction struggling with a
precipitation that is a very stiff function of $W$, and I did not chase it
further.

For the 18 runs that never settle the residual runs from 0.12 to 1690, which is
expected rather than a defect: $\partial\langle h\rangle/\partial t$ is not zero
and the budget as written omits it. Their residual is therefore a useful
independent confirmation of the steadiness classification, not a separate check.

\section{Sensitivity to each parameter in turn}
\label{sec:ladders}

Each ladder is run on both reference states, the positive-stability
$a = 0.79$, $W_c = 40$ ($r = 0.76$) and the negative-stability $a = 0.85$,
$W_c = 50$ ($r = 1.14$). The figures are
\path|05_ladder_dyrad.png| through \path|10_ladder_evap.png|.

\subsection{The radiative contrast $\dyr$: the strongest single lever on the jet}

Raising $\dyr$ from 50 to 110~K takes the jet from 35.8 to 78.3~\ms{} on the
positive-stability reference and from 57.6 to 99.3 on the negative one. The
sensitivity is 0.708 and 0.690~\ms{} per K, nearly the same on both, so this is a
dry-dynamical response that the moisture regime rides on top of rather than
modifies. Peak $|v|$ goes 0.79 to 2.03~\ms{} on the positive reference. The
raining fraction stays 1.000 throughout on the positive side and falls only from
0.648 to 0.591 on the negative side, so forcing strength does not by itself
trigger aggregation over this range.

\textbf{Physical reading (derived).} $\dyr$ sets the meridional temperature
gradient the relaxation drives toward, the pressure gradient in (\ref{eq:v})
follows it, and the jet follows by thermal wind. That the sensitivity is the same
on both sides of the stability boundary says the two effects are close to
separable: latent heating changes where the rain falls and how strong the
overturning is, and $\dyr$ sets the overall amplitude.

\subsection{The eddy moisture diffusivity $D$: the brake, and only on one side}

This is the largest asymmetry in the sweep. On the positive-stability reference,
$D$ from 0 to $4\times10^{6}$~m$^2$\,s$^{-1}$ changes the jet from 53.66 to
53.31~\ms{} (0.7\%) and peak precipitation from 5.77 to 5.61~\mmday{} (3\%). On
the negative-stability reference the same range gives:

\begin{center}
\begin{tabular}{rrrrr}
\toprule
$D$ (m$^2$\,s$^{-1}$) & raining fraction & bands & peak $P$ (\mmday) & $W$ at equator (kg\,m$^{-2}$)\\
\midrule
0                 & 0.046 & 35 & 831.7 & 188.6\\
$2.5\times10^{5}$ & 0.538 & 13 & 212.8 & 85.5\\
$5\times10^{5}$   & 0.568 & 5  & 127.2 & 71.2\\
$1\times10^{6}$   & 0.633 & 5  & 77.8  & 63.0\\
$2\times10^{6}$   & 0.693 & 3  & 52.0  & 58.7\\
$4\times10^{6}$   & 0.718 & 3  & 34.4  & 55.7\\
\bottomrule
\end{tabular}
\end{center}

\textbf{Physical reading (derived).} Section~\ref{sec:budget} showed that the
equilibrium at $D = 10^{6}$ is a near-cancellation between a mean flux importing
$+104.7$~MW\,m$^{-1}$ into the ascending column and an eddy flux exporting
$-106.2$. The eddy flux is proportional to $D$, so at smaller $D$ the ascent has
to build a steeper moisture gradient to remove the same energy, which means a
taller and narrower spike. At $D = 0$ there is nothing to remove it at all, and
the solution collapses to 4.6\% of the domain raining at 832~\mmday{} with the
equatorial column holding 189~kg\,m$^{-2}$ of water. That run also fails the
steadiness test through wandering bands rather than a wandering jet.

Two things follow. The parameter with the most leverage over the aggregated state
is the one the ERA5 calibration is least certain about; its measured range is
$3\times10^{5}$ to $4\times10^{6}$, and over that range the raining fraction goes
0.54 to 0.72. And the positive-stability reference is almost completely
insensitive to it, so $D$ cannot be calibrated from anything on that side.

\begin{figure}[htbp]
\centering
\includegraphics[width=\textwidth]{06_ladder_D.png}
\caption{The eddy moisture diffusivity ladder on both reference states. Hollow
markers are runs that never settle. The positive-stability curves are flat in
every panel; the negative-stability ones are not.}
\label{fig:ladderD}
\end{figure}

\subsection{The convective timescale $\tau_c$: aggregation needs a fast sink}

On the negative-stability reference, $\tau_c$ from 1800 to 172800~s takes the
raining fraction from 0.528 to 0.755 and peak precipitation from 176.8 down to
46.9~\mmday, so a \emph{shorter} convective timescale aggregates more, even
though a shorter $\tau_c$ gives a \emph{smaller} $r$. This is the mechanism
behind the failed collapse of Section~\ref{sec:regime}: concentrating rain
requires the sink to respond quickly to the local moisture, and a slow sink
smears it out. The $\tau_c = 1800$~s run is one of the 18 that never settles.

\subsection{Evaporation $E_0$: the amplitude of everything moist}

$E_0$ barely moves $r$ (1.138 to 1.160 over a fourfold range on the
negative-stability reference) and moves everything else a great deal: peak $|v|$
from 3.31 to 11.15~\ms, peak precipitation from 37.4 to 197.5~\mmday, and the
equatorial temperature from 224.9 to 293.6~K. The temperature is the cleanest
demonstration of the mean-state offset derived in SCIENCE.md~\S3.6: the quiescent
column sits $\tau\Lambda E_0$ above its radiative target, which is 35.6~K at the
lowest $E_0$ and 142.4 at the highest, and over eight runs spanning that range
the measured offset matches the prediction to $4\times10^{-7}$ relative. On the
positive-stability reference $E_0$ changes the jet by
1\% while changing peak precipitation from 3.69 to 9.80~\mmday, exactly tracking
$E_0$ itself because the domain rains uniformly at the evaporation rate.

\subsection{The eddy momentum flux $v_d$: the axisymmetric limit is extreme}

Removing the eddy momentum flux entirely gives a jet of 86.7~\ms{} on the
positive-stability reference and 160.9~\ms{} on the negative one, against 53.6
and 73.8 at the default $v_d = 2.5$~\ms. Peak $|v|$ collapses from 1.27 to
0.31~\ms{} on the positive side, and raising $v_d$ to 7.5 gives 47.5~\ms{} of jet
with 2.07~\ms{} of overturning. All four directions match the expectations in
SCIENCE.md~\S6: stronger eddies give weaker zonal winds, stronger $v$, and a
wider cell.

At $v_d = 0$ the negative-stability run still aggregates, with 76.5\% of the
domain raining and a peak of 72.7~\mmday. So aggregation is thermodynamic, and
the eddy momentum forcing modulates it rather than causing it.

\subsection{The feedback strength $\Lambda$: a threshold, not a gradual onset}

$\Lambda = L_v/C$ is derived rather than free, and detuning it breaks the exact
cancellation of precipitation in the MSE budget, so this ladder is diagnostic
rather than physical. It is also the most informative one. On the
negative-stability reference:

\begin{center}
\begin{tabular}{rrrrr}
\toprule
$\Lambda/\Lambda_0$ & jet (\ms) & peak $|v|$ (\ms) & peak $P$ (\mmday) & raining fraction\\
\midrule
0.00 & 41.58 & 0.562 & 4.91  & 1.000\\
0.25 & 44.69 & 0.683 & 5.16  & 1.000\\
0.50 & 48.86 & 0.912 & 5.70  & 1.000\\
0.75 & 57.36 & 1.691 & 7.93  & 1.000\\
1.00 & 73.79 & 5.860 & 77.82 & 0.633\\
1.25 & 81.90 & 7.526 & 166.85& 0.620\\
\bottomrule
\end{tabular}
\end{center}

Between 0.75 and 1.00 the peak rain rate jumps by a factor of 9.8 and the
raining fraction falls off 1. On the positive-stability reference the same ladder
never aggregates: the raining fraction is 1.000 at every rung and peak
precipitation goes 4.59 to 10.27~\mmday. So the two ingredients are both
necessary: negative gross moist stability, and a feedback strong enough to
exploit it. At $\Lambda/\Lambda_0 = 0$ the two references have identical dry
fields to every digit reported (jet 41.5779~\ms, peak $|v|$ 0.562137~\ms,
$T(0)$ 197.123~K) while their moisture differs as their $W_c$ requires, which is
the expected behaviour of the V2$_0$ bridge and a check that the coupling runs
only in the intended direction.

\begin{figure}[htbp]
\centering
\includegraphics[width=\textwidth]{09_ladder_lambda.png}
\caption{The latent-heating coefficient ladder. The physical value is
$\Lambda_0 = L_v/C = 0.4843$~K per kg\,m$^{-2}$\,s$^{-1}$, the fifth point from
the left. On the negative-stability reference everything jumps between the fourth
and fifth points; on the positive-stability one nothing does.}
\label{fig:ladderlam}
\end{figure}

\subsection{The V1 to V2 attribution chain}

Four runs per reference, changing one thing per step:

\begin{center}
\footnotesize
\begin{tabular}{llrrrrr}
\toprule
$\Hhat$ & step & jet & peak $|v|$ & $T$(0) & $\Delta\theta_{24}$ & peak $P$\\
 & & (\ms) & (\ms) & (K) & (K) & (\mmday)\\
\midrule
either & V1, no heating, no retarget      & 28.20 & 0.365 & 200.3 & 1.13 & 4.38\\
either & retarget to 75 K, $\Lambda=0$    & 41.58 & 0.562 & 197.1 & 1.63 & 4.59\\
\midrule
$+$ & $\Lambda$ on, target still 50 K & 35.84 & 0.790 & 246.9 & 2.38 & 5.08\\
$+$ & $\Lambda$ on, target 75 K (V2)  & 53.57 & 1.266 & 244.6 & 3.33 & 5.72\\
\midrule
$-$ & $\Lambda$ on, target still 50 K & 57.56 & 5.242 & 249.9 & 7.05 & 67.47\\
$-$ & $\Lambda$ on, target 75 K (V2)  & 73.79 & 5.860 & 248.1 & 7.62 & 77.82\\
\bottomrule
\end{tabular}
\end{center}

\noindent The first two rows have no moisture feedback, so they are the same run
whichever $W_c$ is set; the peak $P$ differs only because $P$ reads off $W_c$.

The retarget alone accounts for a 47\% jet increase with no moisture feedback at
all, which is where most of the overshoot against observations comes from. The
71~K warming appears with $\Lambda$, as derived. And aggregation appears in the
negative-stability row as soon as $\Lambda$ is switched on, even at the flatter
50~K target, so it does not require the steeper radiative contrast.

\section{Two equilibria at the same parameters?}
\label{sec:multistable}

A state whose ascent reinforces itself is a natural candidate for multiple
equilibria, so eight runs start with the column already moist: at
$a = 0.85$ and $W_c = 40$, 42, 44 and 46, spanning $r$ from 0.919 to 1.055
across the transition, with $W(t=0) = W_c + 15$ and $W_c + 30$~kg\,m$^{-2}$
instead of the default $W_c$.

Nothing happens. The largest relative difference from the default start, over any
of the reported quantities and any of the eight runs, is
$8\times10^{-4}$. At $W_c = 40$, 42 and 44 the equilibria agree to every digit
reported: jet 60.81, 64.68, 70.50~\ms{} and peak $P$ 9.02, 15.02,
38.04~\mmday{} whatever the initial moisture. At $W_c = 46$ the jet differs by
0.06~\ms{} out of 74.07 and peak precipitation by 0.04 out of 56.64.

So within the range tested there is one attractor per parameter set, and the
aggregated state is reached from a dry start as readily as from a wet one. That
is a negative result and worth having: it means the regime map of
Section~\ref{sec:regime} is a map of the parameter space rather than of the
initial conditions, and the boundary at $r \approx 0.94$ is not a hysteresis
loop. What it does not test is a start that is drier than $W_c$, or an
asymmetric initial perturbation, either of which could still find a second
branch (Section~\ref{sec:open}).

\section{Forcing peaked off the equator and held fixed in time}
\label{sec:offeq}

Sixteen runs move the maximum of the radiative target off the equator and hold
it there, at $y_0 = 0$, 700, 1400, 2100 and 2800~km. Two profile families are
used. The $\sin^2$ family translates rigidly, so $y_0$ displaces the forcing
without changing its shape or strength. The SB08 family adds an antisymmetric
term, so its peak value \emph{rises} above $\theta_{00}$ by
$\dyr\sin^2(\pi y_0/2y_1)$ as $y_0$ increases: a solstitial SB08 forcing is
stronger as well as displaced, which is worth holding in mind when its
circulation is stronger too. Both families are run on both reference states.

\subsection{Where the rain goes, and how to ask the question}

\textbf{Three positions, and they are not interchangeable.} The precipitation
centroid taken over the whole domain is a bad tracker here, by enough to change
the answer. On the positive-stability reference the domain drizzles at the
evaporation rate everywhere, and that uniform background pulls the centroid
toward the middle of the domain: at $y_0 = 1400$~km the plain centroid sits at
0.299~Mm while the rain \emph{maximum} sits at 1.082~Mm. Read from the centroid
the rain band follows 21\% of the forcing displacement, read from the maximum
77\%. The report uses the maximum, and a centroid taken over the connected
region raining above $E_0$, which agree with each other. All three are in
\path|all_runs.csv|, so the difference is visible rather than buried.

\textbf{The rain follows most but not all of the forcing.} On the
positive-stability reference the rain maximum reaches 1.082, 1.404 and
1.711~Mm for SB08 forcing at $y_0 = 1400$, 2100 and 2800~km, so it covers 77\%,
67\% and 61\% of the displacement. The fraction falls as the forcing goes
further out, which is what a restoring influence tied to the equator produces.
Under $\sin^2$ forcing the fractions are slightly higher, 82\% at 1400~km and
70\% at 2100~km.

\subsection{The rain maximum and the energy flux equator are 1.9 Mm apart}

Take the SB08 $y_0 = 1400$~km run on the positive-stability reference. Its rain
peaks at 1.082~Mm. The zero of its total column MSE flux, the energy flux
equator, is at 2.959~Mm. The two are separated by 1.88~Mm, which is about
$15^\circ$ of equivalent latitude.

The reason is visible in the flux itself. The gross moist stability in that run
varies only between 17.8 and 19.3~MJ\,m$^{-2}$ across the whole tropics, and the
eddy flux never exceeds 0.8~MW\,m$^{-1}$ against a mean flux reaching
80~MW\,m$^{-1}$. So the total flux is very nearly $v\Hhat$ with $\Hhat$
constant, which means it vanishes where $v$ does. And $v$ has no zero anywhere
between $-5.0$ and $+3.0$~Mm: a single cross-equatorial cell spans the entire
tropics with southward upper-level flow throughout, and the rain maximum sits
\emph{inside} that cell rather than at its edge. Upper-level divergence
$\py v$ is positive across the whole band, peaking at 1.061~Mm, which is where
the rain is (1.082~Mm, agreeing to 21~km). The flux zero at 2.959~Mm is the
boundary between that cell and the weak summer cell poleward of it.

So in this model the standard reasoning that puts the rain band at the energy
flux equator does not apply, and the reason is structural: the ascending region
is broad and lies within a single-signed $v$, so the energy flux does not have
to vanish where the ascent is strongest. Whether that is a defect of the
single-layer formulation or a genuine consequence of a nearly uniform $\Hhat$
is the sharpest open question this sweep raises
(Section~\ref{sec:open}).

\subsection{The winter jet reaches four times the observed value}

Under SB08 forcing on the positive-stability reference, as $y_0$ goes from 1400
to 2800~km:

\begin{itemize}
\item The winter (southern) jet strengthens from 91.6 to 128.8~\ms{} while the
summer jet weakens from 26.9 to 12.3~\ms.
\item Equatorial zonal wind becomes strongly easterly, from $-25.4$ to
$-54.2$~\ms.
\item Peak $|v|$ more than doubles, 4.38 to 9.50~\ms.
\item The domain stops raining everywhere: the raining fraction falls from
1.000 to 0.845, so even a positive-stability state develops a dry zone once the
forcing is far enough off the equator.
\end{itemize}

The directions are all as expected for a solstitial circulation, and the
magnitudes are not. A 129~\ms{} winter jet is four times the observed
zonal-mean maximum of 31~\ms, and $-54$~\ms{} of equatorial easterly is roughly
three times what the solstice season shows. Section~\ref{sec:plausible} takes
the plausibility question up properly; the point here is that the off-equatorial
overshoot compounds the jet overshoot the equinoctial runs already have.

\subsection{On the negative-stability reference the rain leaves the tropics}

The same forcing on the negative-stability reference ($a = 0.85$, $W_c = 50$)
does something qualitatively different. At $y_0 = 700$~km the rain maximum is
at 3.44~Mm, and by $y_0 = 2800$~km it is at 5.28~Mm, which is $46^\circ$
equivalent latitude. The rain band is no longer near the forcing maximum at all;
it has been relocated far poleward, and it now sits \emph{poleward} of the
energy flux equator ($-0.01$ to $2.48$~Mm) rather than equatorward of it.

Every one of these eight runs fails the steadiness test of
Section~\ref{sec:numerics}, so the positions above are maxima of a 1000-day mean
field rather than of a settled state, and they should be read as such. The
cleanest statement the block makes is that split itself: with the forcing off
the equator, every run on the negative-stability reference keeps fluctuating and
every run on the positive-stability reference settles.

\begin{figure}[htbp]
\centering
\includegraphics[width=\textwidth]{11_offequatorial.png}
\caption{Time-invariant off-equatorial forcing. Rain position, energy flux
equator, equatorial zonal wind, the two jets and the overturning strength against
the forcing displacement $y_0$, for both profile families and both reference
states. The dotted 1:1 line in the first panel marks where the rain would sit if
it followed the forcing exactly.}
\label{fig:offeq}
\end{figure}

\begin{figure}[htbp]
\centering
\includegraphics[width=\textwidth]{12_offeq_profiles_P.png}
\caption{Equilibrium profiles for the SB08 off-equatorial series on the
positive-stability reference, $y_0$ from 0 to 2800~km. The winter jet near
$y = -5$~Mm strengthens to 129~\ms{} while the summer jet nearly disappears.}
\label{fig:offeqprof}
\end{figure}

\section{Forcing that migrates through a seasonal cycle}
\label{sec:seasonal}

Nineteen runs vary $y_0$ in time with the SB08 profile: five amplitudes from 350
to 2800~km at a 360-day period, five periods from 90 to 1440 days at 1400~km
amplitude, and the square and tanh waveforms, each on both reference states, plus
one $n_y = 1601$ twin. Each is integrated for 2880~days plus at least six full
cycles, and the last two complete cycles are compared as the repeatability
check. All amplitudes and lags come from a least-squares fit at the forcing
frequency, which is more robust on these series than a peak-finding estimate.

\begin{center}
{\footnotesize
\begin{tabular}{lcccccccc}
\toprule
run & amp. & period & $\hat H$ & centroid amp. & gain & lag & summer jet & winter jet\\
 & (km) & (d) & & (km) & & (d) & (m\,s$^{-1}$) & (m\,s$^{-1}$)\\
\midrule
\texttt{M\_P\_amp1400\_p0090} & 1400 & 90 & $+$ & 1192 & 0.852 & 3.1 & 70.2 & 70.2\\
\texttt{M\_P\_amp1400\_p0180} & 1400 & 180 & $+$ & 1254 & 0.896 & 4.1 & 82.1 & 82.1\\
\texttt{M\_P\_amp0350\_p360} & 350 & 360 & $+$ & 390 & 1.115 & 3.9 & 61.3 & 61.3\\
\texttt{M\_P\_amp0700\_p360} & 700 & 360 & $+$ & 753 & 1.076 & 4.3 & 70.0 & 70.0\\
\texttt{M\_P\_amp1400\_p360} & 1400 & 360 & $+$ & 1288 & 0.920 & 5.3 & 89.1 & 89.1\\
\texttt{M\_P\_amp1400\_p360\_tanh} & 1400 & 360 & $+$ & 1265 & 0.906 & 5.1 & 91.0 & 91.0\\
\texttt{M\_P\_amp1400\_p360\_square} & 1400 & 360 & $+$ & 1105 & 0.972 & 3.8 & 91.4 & 91.4\\
\texttt{O\_P\_amp1400\_p360\_ny1601} & 1400 & 360 & $+$ & 1287 & 0.920 & 5.3 & 89.1 & 89.1\\
\texttt{M\_P\_amp2100\_p360} & 2100 & 360 & $+$ & 1626 & 0.774 & 6.0 & 108.3 & 108.3\\
\texttt{M\_P\_amp2800\_p360} & 2800 & 360 & $+$ & 1957 & 0.699 & 6.2 & 126.7 & 126.7\\
\texttt{M\_P\_amp1400\_p0720} & 1400 & 720 & $+$ & 1311 & 0.937 & 5.9 & 91.1 & 91.1\\
\texttt{M\_P\_amp1400\_p1440} & 1400 & 1440 & $+$ & 1321 & 0.944 & 6.0 & 91.4 & 91.4\\
\texttt{M\_N\_amp1400\_p0090} & 1400 & 90 & $-$ & 4116 & 2.941 & 4.8 & 122.6 & 140.2\\
\texttt{M\_N\_amp0350\_p360} & 350 & 360 & $-$ & 384 & 1.097 & -162.4 & 120.1 & 110.3\\
\texttt{M\_N\_amp0700\_p360} & 700 & 360 & $-$ & 2759 & 3.942 & 42.5 & 133.7 & 137.3\\
\texttt{M\_N\_amp1400\_p360} & 1400 & 360 & $-$ & 3132 & 2.237 & 14.9 & 166.3 & 170.2\\
\texttt{M\_N\_amp2100\_p360} & 2100 & 360 & $-$ & 3634 & 1.731 & 14.2 & 196.6 & 193.0\\
\texttt{M\_N\_amp2800\_p360} & 2800 & 360 & $-$ & 5661 & 2.022 & 1.0 & 200.1 & 217.1\\
\texttt{M\_N\_amp1400\_p1440} & 1400 & 1440 & $-$ & 2630 & 1.879 & 135.3 & 176.9 & 178.5\\
\bottomrule
\end{tabular}
}

\end{center}

\subsection{The positive-stability cycles repeat exactly, and the negative ones never repeat}

The cleanest result of the block is the split. On the positive-stability
reference, consecutive cycles differ by 0.0~km in rain-band position for nine of
the ten runs, and 0.4~km for the square-wave one. That is a fully converged
periodic solution. On the negative-stability reference every run differs from its
own previous cycle by 10,889 to 14,656~km, which is most of the domain. Those are
chaotic, and their fitted amplitudes and lags describe one arbitrary realisation.
They are in the table and marked, and no conclusion rests on them.

\subsection{The response is quasi-static, and the obvious null model is wrong by a factor of ten}

\begin{figure}[htbp]
\centering
\includegraphics[width=\textwidth]{16b_seasonal_period.png}
\caption{Response against forcing period at fixed 1400~km amplitude, with the
prediction of a single relaxation on the model's thermal timescale $\tau = 37$~d
in grey. Dashed lines with open markers are the chaotic negative-stability runs,
whose fitted values describe one realisation.}
\label{fig:seasperiod}
\end{figure}

A single first-order relaxation on the model's own thermal timescale
$\tau = 37$~days is the natural null model. Driving $dX/dt = (F(t)-X)/\tau$ with
$F = A\sin\omega t$ gives a response damped by $1/\sqrt{1+\omega^2\tau^2}$ and
lagged by $\arctan(\omega\tau)/\omega$. That predicts a lag of 17, 26, 33, 36 and
37~days at periods of 90, 180, 360, 720 and 1440~days, and a damping factor of
0.36, 0.61, 0.84, 0.95 and 0.99.

What the runs give at those five periods is a lag of 3.1, 4.1, 5.3, 5.9 and
6.0~days, and a gain of 0.852, 0.896, 0.920, 0.937 and 0.944. The lag is five to
six times smaller than predicted at every period, and the damping at the shortest
period is 0.85 where the null model says 0.36.

\textbf{Physical reading (derived, then checked).} The prediction fails because
the rain band's position is not something the column temperature has to
equilibrate to. It is set by where the forcing peaks, moderated by whatever pulls
the ascent back toward the equator, and both act on the fast dynamical and
moisture timescales rather than on $\tau$. The check is that the seasonal gain
should then equal the \emph{steady} response at the same displacement. It does:
the time-invariant SB08 run at $y_0 = 1400$~km puts the rain maximum at
1.082~Mm, a gain of 0.773, and the seasonal runs at 1400~km amplitude give a
rain-maximum gain of 0.757 to 0.787 across the five periods, so the steady value
sits inside that range. The migrating solution is tracking the sequence of steady
states.

The same holds for the dynamics. Local Rossby number in the dominant cell,
averaged over $|y|$ from 1 to 4~Mm, reaches 0.452, 0.540, 0.597, 0.616 and 0.622
at the five periods, and the steady $y_0 = 1400$~km run gives 0.624. At the two
longest periods the migrating cycle matches the perpetual state to 1\%, and at a
360-day period it falls 4\% short. So there is no sign in this model of the
seasonal cycle reaching an angular-momentum-conserving regime that the
corresponding steady forcing does not, which is the mechanism Schneider and
Bordoni identify for monsoon onset. If anything the migration keeps the cell
slightly further from that limit.

\subsection{What the amplitude does}

\begin{figure}[htbp]
\centering
\includegraphics[width=\textwidth]{13_seasonal_hov_P.png}
\caption{Precipitation through the composite cycle at five migration amplitudes
on the positive-stability reference. Four position measures are drawn: the
forcing maximum, the rain maximum, the centroid of the rain band, and the
maximum of upper-level divergence. On the positive-stability reference they
coincide closely.}
\label{fig:seashov}
\end{figure}

Raising the migration amplitude from 350 to 2800~km at a 360-day period takes
the summer-side jet from 61.3 to 126.7~\ms{} and the rain-band amplitude from 390
to 1957~km, so the gain falls from 1.115 to 0.699. The gain exceeding one at the
smallest amplitudes is worth flagging as a small-signal result rather than
interpreting: at 350~km amplitude the band moves 390~km, and the difference is
of the order of ten gridpoints.

At amplitudes of 2100 and 2800~km a dry region opens in the winter hemisphere
even on the positive-stability reference. The raining fraction dips to 0.905 and
0.844 at the solstices of the migrating runs, and the steady runs held at the
same displacement give 0.905 and 0.845. So a far enough displacement produces a
dry zone with no help from negative gross moist stability, and the migrating and
steady cases produce the same one.

\subsection{A migrating forcing against one held at the solstice}

Four pairs compare the seasonal solution at its peak $y_0$ against the
time-invariant run with the forcing held at that same $y_0$:

\begin{center}
\begin{tabular}{lrrrr}
\toprule
 & rain position, seasonal & steady & summer jet, seasonal & steady\\
\midrule
700~km, positive $\Hhat$  & 0.102~Mm & 0.124~Mm & 44.7~\ms & 38.4~\ms\\
1400~km, positive $\Hhat$ & 0.262    & 0.299    & 38.6     & 26.9\\
2100~km, positive $\Hhat$ & 0.476    & 0.529    & 33.8     & 18.7\\
\bottomrule
\end{tabular}
\end{center}

The migrating case reaches a rain band 22 to 53~km closer to the equator and a
summer jet 6.3 to 15.1~\ms{} stronger than the state it would reach if the
forcing stopped there. Both differences point the same way and have the same
cause: the summer hemisphere has not had time to finish weakening its jet or to
finish pulling the rain poleward. The peak precipitation is essentially identical
between the two (6.02 against 6.03, 6.86 against 6.88, 8.02 against
8.06~\mmday), so the migration changes the geometry of the circulation without
changing its rainfall intensity.

The fourth pair, at 1400~km on the negative-stability reference, differs far more
(rain at 2.509 against 1.408~Mm, jet 113.8 against 82.2~\ms, peak precipitation
229.4 against 50.6~\mmday), and both members of that pair fail their respective
steadiness and repeatability tests, so the comparison is between two fluctuating
states and is not interpretable.

\subsection{Waveform}

Replacing the sinusoidal migration with a tanh or a square wave at 1400~km and
360~days changes very little: the gain is 0.920, 0.906 and 0.972 for sin, tanh
and square, the lag 5.3, 5.1 and 3.8~days, and the summer jet 89.1, 91.0 and
91.4~\ms. The square wave's smaller lag is the expected consequence of a forcing
that spends most of the cycle stationary. The square-wave run is also the only
positive-stability seasonal run whose consecutive cycles are not identical, and
its 0.4~km difference is small enough to be the discrete jump itself rather than
a failure to converge.

\section{Timestep, grid spacing, and how far from equilibrium the runs are}
\label{sec:numerics}

\subsection{The documented timestep is not enough everywhere}

Before spending hours on long integrations I ran all 149 configurations for 100
days each and checked the output files. Twenty-two of them diverged, meaning the
integration hit a NaN and stopped, in most cases within the first day or two.
The repository's note that $\Delta t = 30$~s is the V2 ceiling at $n_y = 801$ is
right at the reference parameters and wrong as a general statement.

The 22 fall into two groups, and both make sense. One group has an unusually
strong circulation: the three deeply negative-stability cases at $W_c = 60$ with
$a \ge 0.83$ ($r$ from 1.29 to 1.45), the two steepest radiative contrasts
$\dyr = 100$ and $110$~K on the negative-stability reference, the two smallest
static stabilities $\Delta_z = 30$ and $45$~K, and the one case with the latent
heating coefficient raised 25\% above its physical value. The other group has a
forcing maximum far off the equator: every run with $|y_0| \ge 1400$~km, and the
$y_0 = 700$~km runs on the negative-stability reference.

Halving to $\Delta t = 15$~s rescued 16 of the 22 over a 300-day probe. The
remaining six needed $\Delta t = 6$~s, and all five that stayed in the sweep ran
clean at that value; the sixth was an $n_y = 1601$ off-equatorial case that would
have cost 2.4~h on its own, so it was replaced (Section~\ref{sec:offeq}).

\subsection{These are cold-start transients, not unstable equilibria}

The distinction matters, because it says the equilibria themselves are
reachable. A run whose forcing maximum sits at $y_0 = 1400$~km from the first
timestep diverges at $\Delta t = 30$~s. A \emph{seasonal} run whose forcing
maximum walks out to the same 1400~km over 90 days does not: it survived the
100-day screen at $\Delta t = 30$~s. The model starts from $u = v = 0$ with
$\theta$ at its target, and a large off-equatorial forcing imposed on that state
produces a violent adjustment; the same forcing approached gradually does not.

Even so, every seasonal run in the sweep uses $\Delta t = 15$~s. Two reasons.
One seasonal case, amplitude 2800~km on the negative-stability reference,
diverged at day 800 rather than in the first days, so the screen would have
passed it and the production run would have wasted its time. And holding
$\Delta t$ fixed across the seasonal block removes the timestep as a confounder
when amplitudes and periods are compared against each other.

\subsection{Timestep and grid spacing}

Eight twin pairs test the numerics, each changing only $\Delta t$ or $n_y$ from
its partner. The comparison is the pointwise difference after interpolating to
the coarser grid, not a difference of summary numbers.

\begin{center}
\small
\begin{tabular}{lrrrr}
\toprule
pair & max$|\Delta u|$ (\ms) & as \% of peak $u$ & where (Mm) & jet change\\
\midrule
positive $\Hhat$: $\Delta t$ 30 to 15 & 0.0031 & 0.01\% & $-11.5$ & 0.00\%\\
negative $\Hhat$: $\Delta t$ 30 to 15 & 0.0087 & 0.01\% & $-7.4$  & 0.00\%\\
positive $\Hhat$: $n_y$ 801 to 1601   & 0.589  & 1.10\% & $-9.7$  & 0.02\%\\
negative $\Hhat$: $n_y$ 801 to 1601   & 7.20   & 9.76\% & $-7.4$  & 0.04\%\\
near-critical: $n_y$ 801 to 1601      & 0.260  & 0.37\% & $-11.1$ & 0.02\%\\
off-equatorial: $n_y$ 801 to 1601     & 29.98  & 41.58\%& $-7.5$  & 0.02\%\\
positive $\Hhat$: $n_y$ 801 to 401    & 7.58   & 14.14\%& $-9.6$  & 0.01\%\\
negative $\Hhat$: $n_y$ 801 to 401    & 49.86  & 67.58\%& $+7.4$  & 0.30\%\\
\bottomrule
\end{tabular}
\end{center}

Two conclusions, and they point different ways.

\textbf{The timestep is converged and the bulk quantities are converged in
resolution.} Halving $\Delta t$ changes $u$ nowhere by more than 0.01\% of its
peak, which is negligible; that matters because 21 of the 151 runs use a smaller
$\Delta t$ than the rest for stability, and this says mixing timesteps across the
sweep costs nothing. Across all eight pairs the jet moves by at most 0.30\% and
peak precipitation by at most 0.91\%, so every scalar this report quotes is
resolution-insensitive.

\textbf{The profile in the westerly terminus region is not converged.} Every one
of the large differences sits at $|y|$ between 7 and 11~Mm, which is where the
zonal wind turns over and the Heaviside gate on the eddy momentum flux switches.
The 2026-07-12 $n_y = 1601$ probe recorded the same thing for the dry model, so
this is a known property of the formulation rather than something the moisture
introduces. What is new is the size on the negative-stability side: 9.8\% between
801 and 1601, and 67.6\% between 801 and 401. Any conclusion about the
extratropical structure of an aggregated solution should be treated as
resolution-dependent. Nothing in this report rests on one.

The off-equatorial pair is the weakest of the eight, and deliberately so: the
$y_0 = 700$~km cold start diverges at $n_y = 1601$ unless $\Delta t = 6$~s, so
both members of that pair run at 6~s and only for 2500 days, which is 2.2 drag
timescales rather than 5. Both are equally unequilibrated, so their difference
still measures grid sensitivity, and their bulk agreement is good (jet 38.44
against 38.45~\ms, peak $|v|$ 2.424 against 2.426).

\subsection{How far from equilibrium the runs are}

Every non-seasonal run is a fixed 5800-day integration, which is 5.01 drag
timescales at $\epsilon_u = 10^{-8}$~s$^{-1}$, rather than a convergence-gated
one. That choice makes the drift something this report measures instead of
something a detector asserts.

For the 113 runs classified steady, the pointwise difference in $u$ between the
mean over days 3800--4800 and the mean over days 4800--5800 is between
$7\times10^{-10}$ and $3.9\times10^{-3}$ of peak $u$, and the two runs at the top
of that range are the 2500-day off-equatorial pair. Excluding those, the largest
drift among the 111 full-length steady runs is $4.5\times10^{-4}$. The domain
integral of
precipitation equals that of evaporation to between 0.999 and 1.000 in the
finite-volume measure for all 131, which is a constraint the integration has to
satisfy at equilibrium and does. And no run had negative $W$ at any timestep,
which the MUSCL-limited transport is designed to prevent.

\subsection{What the never-settling solutions look like}

The 18 unsteady runs fall into two kinds. In 16 of them the jet magnitude itself
swings, by 2.6 to 83.0~\ms{} over the trailing 2000 days. In two, both from the
low-$D$ end of the diffusivity ladder, the jet holds to under 0.7~\ms{} while the
rain bands wander: their pointwise $u$ difference between consecutive 1000-day
means is 10\% and 62\%. A steadiness test on the jet alone would have passed
those two, and their MSE budget residuals of 0.12 and 0.45 of the source are what
gave them away.

\begin{figure}[htbp]
\centering
\includegraphics[width=\textwidth]{18_spinup.png}
\caption{Equilibration of five runs spanning the transition, against the 1157-day
drag timescale. The $r = 1.37$ run is one of the 18 that never settles; its jet
and peak rain rate keep swinging for the whole integration.}
\label{fig:spinup}
\end{figure}

The unsteadiness shows up in the profiles as large-amplitude oscillations in $u$
poleward of about 5~Mm, visible on the two largest-$W_c$ curves of
Figure~\ref{fig:transect} reaching $\pm80$~\ms. Their wavelength is roughly
1.6~Mm, which is 40 gridpoints at $n_y = 801$, so they are not the $2\Delta y$
computational mode. Whether they are resolution-converged is untested: no
$n_y = 1601$ twin was run above $r = 1.01$, and given the 9.8\% difference the
twins do show at $r = 1.145$ in exactly that region, they should not be trusted
quantitatively.

\section{Open questions, and the runs that would settle them}
\label{sec:open}

Each of these names what outcome would decide it, because a proposed run without
that is not a test.

\paragraph{1. Why is the rain band 1.9 Mm from the energy flux equator, and is
that a defect of the single layer?} This is the sharpest question the sweep
raises, because the standard framework for ITCZ position ties the two together
and here they come apart. The argument in Section~\ref{sec:offeq} is that with
$\Hhat$ nearly uniform the total flux is essentially $v\Hhat$, so it vanishes
where $v$ does, which is the cell boundary. That argument is testable: it predicts
that the separation should shrink as $\Hhat$ becomes less uniform. The
negative-stability off-equatorial runs have strongly varying $\Hhat$ and are all
unsteady, so they cannot answer it. A run at intermediate $r$, around 1.01 to
1.05, with off-equatorial forcing, would have varying $\Hhat$ and might still
settle. If the separation shrinks with the variance of $\Hhat$ the explanation
holds; if it does not, something else is going on. This is the run I would do
first.

\paragraph{2. Is the 40-gridpoint extratropical wave train real?} The
never-settling solutions oscillate at a wavelength of about 1.6~Mm poleward of
5~Mm. It is not the grid-scale computational mode, and the twins at $r = 1.145$
already show a 9.8\% resolution sensitivity in exactly that region. An
$n_y = 1601$ twin of one unsteady run, say $a = 0.85$, $W_c = 60$, would say
whether the oscillation survives grid refinement. At $\Delta t = 15$~s and 5800
days that costs about an hour. If the wave train changes character with
resolution, then the whole $r > 1.3$ corner of the map is describing the numerics
rather than the physics.

\paragraph{3. Is there a second branch reachable from a drier start?} The
initial-condition block only tested moister starts, and found one attractor.
Starting below $W_c$, so that nothing precipitates at $t = 0$, tests the other
direction, and so does an asymmetric initial perturbation, which could find a
single-cell state. Eight runs, about two hours. A negative result there would
make the regime map genuinely a map of parameter space.

\paragraph{4. What is $D$, really?} $D$ has more leverage over the aggregated
state than any other parameter, and the ERA5 calibration puts it only within a
factor of 13, from $3\times10^{5}$ to $4\times10^{6}$~m$^2$\,s$^{-1}$. Across that
range the raining fraction of the negative-stability reference goes 0.54 to 0.72
and its peak rain rate 213 to 34~\mmday. No model result on that side of the
boundary is better determined than $D$ is. The ERA5 report's own recommendation,
downloading the five empty ERA5 flux directories it lists, is the way to narrow
it, and this sweep raises the value of doing so.

\paragraph{5. Can the jet overshoot be fixed without breaking the retarget?} The
jet is 1.5 to 2.6 times observed, and Section~\ref{sec:plausible} traces most of
it to the radiative retarget: 0.708~\ms{} per K of $\dyr$, needing
$\dyr \approx 43$~K to reach the observed 31~\ms, which is below the
radiative-convective 50~K and therefore not available. Two things are worth
trying and neither was in this sweep. A larger $v_d$ reduces the jet without
touching the thermodynamics: at $v_d = 7.5$~\ms{} the jet is 47.5~\ms{} on the
positive-stability reference, still 1.5 times observed, so a $v_d$ ladder reaching
15 or 20 would say whether eddy momentum forcing alone can close the gap and at
what cost to the overturning. And making the radiative background \emph{colder}
as well as steeper is the refinement SCIENCE.md~\S3.6 flags as not yet done; it is
dynamically free as implemented, so it would not help the jet, which is itself
worth stating so nobody tries it expecting otherwise.

\paragraph{6. Does the eddy energy transport gap have a fix inside this
formulation?} The model transports 0.38~MW\,m$^{-1}$ of eddy energy at
$24^\circ$ against an observed 57.7, and the reason is structural: there is no
eddy dry-static-energy flux in the equations, and the one eddy flux there is acts
on a $W$ field that local evaporation and precipitation hold nearly flat. Adding
$\kappa_\theta\,\py^2\theta$ is already in the code and disabled by default. A
ladder in $\kappa_\theta$ would show whether an eddy heat flux large enough to
matter for the energy budget leaves the circulation recognisable. That is a small
change to test and would answer whether the gap is a missing term or a missing
layer.

\paragraph{7. Which value of $r$ should the model actually run at?} The
observations put the tropics near $r = 1$: the ERA5 estimate of $\Wstar$, 44 to
57~kg\,m$^{-2}$, straddles the model's quiescent 50.66, and the observed
band-mean $\Hhat/\Shat$ of 0.26 to 0.32 is positive but small. This sweep says
that the model's behaviour on the two sides of $r = 1$ differs in kind, and that
the boundary sits at $r \approx 0.94$ rather than 1. So the choice of $a$ is not
a detail. The ERA5-recommended $a \approx 0.79$ with $W_c = 40$ puts the model at
$r = 0.76$, comfortably on the smooth side, and its jet is the least unrealistic
of the configurations tested. Whether the real tropical atmosphere sits far
enough onto the smooth side to justify that is a question about the observations,
not about the model, and the largest uncertainty in it is the 34\% floor on
$\Shat$ from the barotropic mass adjustment that the ERA5 report documents.

\appendix

\section{The model equations, as the code implements them}
\label{sec:equations}

Transcribed from \path|ss09/sw_model.py| and \path|ss09/sw_config.py|, and
cross-checked against \path|SCIENCE.md| \S2--3. Not written from memory: the
function that computes each term is named beside it, because an earlier report
in this repository had four of these equations wrong.

The model is a single layer representing the zonal-mean upper troposphere on an
equatorial $\beta$-plane, $y \in [-15751, +15751]$~km with rigid walls. The
prognostic variables are the zonal wind $u$ and potential temperature $\theta$
on $n_y$ cell centres, the meridional wind $v$ on the $n_y-1$ interior faces
(Arakawa C-grid), and the column water vapour $W$ on the centres.

\paragraph{Zonal momentum} (\texttt{StaggeredSWModel.du\_dt})
\begin{equation}
\frac{\partial u}{\partial t}
= \underbrace{\beta y\, \bar{v}}_{\texttt{coriolis\_term}}
- \underbrace{\bar{v}\,\py u\big|_{\rm upwind}}_{\texttt{merid\_advec\_u}}
- \underbrace{u\,(\py v)\,\mathcal{H}(\py v)}_{\texttt{vert\_advec\_u}}
- \underbrace{\epsilon_u u}_{\texttt{rayleigh\_drag\_u}}
- \underbrace{v_d\,\mathcal{H}(u)\,\mathrm{sgn}(y)\,\py u}_{\texttt{edd\_mom\_flux\_div\_u}} .
\label{eq:u}
\end{equation}
$\bar v$ is the face $v$ averaged to the centres (\texttt{v\_faces\_to\_centers});
$\py v$ is the compact two-point face difference
(\texttt{v\_divergence\_at\_centers}). The eddy momentum flux divergence is
\emph{advective} in form, transporting $u$ poleward at speed $v_d$ wherever the
flow is westerly, and its $\py u$ uses a MUSCL reconstruction with
monotonized-central limited slopes (\texttt{muscl\_mc\_du\_dy}), which is the
production stencil. Note the gate on the vertical momentum advection is the
kinematic $\mathcal{H}(\py v)$, SS09's own form, not the thermodynamic
$\mathcal{H}(\theta_E-\theta)$ used for dry and V1 runs: under latent heating a
precipitating column runs \emph{warmer} than its relaxation target, so the
thermodynamic gate would switch off exactly where convection is strongest.

\paragraph{Meridional momentum} (\texttt{StaggeredSWModel.dv\_dt}), on the faces
\begin{equation}
2\frac{\partial v}{\partial t}
= -\beta y\,\bar{u}
- \frac{gH}{T_0}\,\py T
+ k_v\,\py^2 v ,
\qquad T = \theta\,(p_t/p_s)^{R/c_p} = \theta/1.6 .
\label{eq:v}
\end{equation}
$\bar u$ is $u$ averaged to the faces, $\py T$ the compact face difference. The
factor of 2 is the vertical integration with $v_{\rm top} = -v_{\rm bottom}$.
The Laplacian uses a mirror wall ghost so that $v = 0$ at the walls.

\paragraph{Thermodynamics} (\texttt{dtheta\_dt}, \texttt{latent\_heating})
\begin{equation}
\frac{\partial \theta}{\partial t}
= \frac{\theta_{\rm rad} - \theta}{\tau}
- \frac{d\,\Delta_z}{H}\,\py v
+ \Lambda P ,
\qquad
\Lambda = \frac{L_v}{C},
\quad
C = \frac{c_p\,\Pi\,\Delta p}{g} .
\label{eq:theta}
\end{equation}
This equation is \emph{not} gated. $\theta_{\rm rad}$ is the same forcing-profile
family as $\theta_E$ but at the steeper purely radiative contrast $\dyr$
(\texttt{radiative\_target\_profile}), because the convective heating that V1's
single relaxation bundled in is now explicit as $\Lambda P$ and keeping both
would double-count it.

\paragraph{Column water vapour} (\texttt{moisture\_transport\_tendency},
\texttt{precipitation})
\begin{equation}
\frac{\partial W}{\partial t}
+ \py\Big[-(2a-1)\,v\,W - D\,\py W\Big]
= E_0 - P ,
\qquad
P = \frac{(W - W_c)^{+}}{\tau_c} .
\label{eq:w}
\end{equation}
$-(2a-1)v$ is the column transport velocity: the lower branch moves opposite to
the upper-layer $v$, and with a fraction $a$ of the column water vapour in the
lower branch the column-integrated flux reduces to this one coefficient. The
transport is finite-volume in flux form through the faces, with MUSCL-MC face
values of $W$ upwinded on the transport velocity's sign and zero flux at the
walls, so transport moves no total water to roundoff. Advection is stepped at
the central leapfrog level; $D$ and $P$ at the lagged $n-1$ level.

\paragraph{Time stepping} Leapfrog with an Asselin filter, coefficient 0.04.
$\theta$ and $W$ share the same $\Delta t$ and filter coefficient, and the
latent heating reads the \emph{same} lagged $P$ array the moisture sink uses,
so $C\Lambda = L_v$ holds discretely and precipitation cancels exactly between
the two budgets.

\paragraph{Column energetics} Combining (\ref{eq:theta}) and (\ref{eq:w}) with
$\langle h\rangle = C\theta + L_v W$,
\begin{equation}
\frac{\partial \langle h\rangle}{\partial t}
+ \py\Big[v\Hhat - L_v D\,\py W\Big]
= \frac{C}{\tau}\big(\theta_{\rm rad} - \theta\big) + L_v E_0 ,
\qquad
\Hhat = \Shat - L_v(2a-1)W ,
\label{eq:mse}
\end{equation}
with the gross dry stability $\Shat = C\,d\,\Delta_z/H$. Precipitation is
internal and does not appear.

\paragraph{Forcing profiles} The $\sin^2$ profile
(\texttt{Sin2Profile.\_\_call\_\_}) is
\begin{equation}
\theta_E = \theta_{00} - \Delta_y \sin^2\!\left(\frac{\pi (y - y_0)}{2 y_1}\right),
\qquad |y - y_0| < y_1 ,
\end{equation}
a rigid translation of the equinoctial profile. The SB08 profile
(\texttt{SB08Profile.\_\_call\_\_}) is
\begin{equation}
\theta_E = \theta_{00} - \Delta_y\left[
\sin^2\!\left(\frac{\pi y}{2 y_1}\right)
- 2 \sin\!\left(\frac{\pi y_0}{2 y_1}\right)\sin\!\left(\frac{\pi y}{2 y_1}\right)
\right],
\end{equation}
clamped to its $\pm y_1$ values outside. Its maximum sits at $y = y_0$, and its
peak value \emph{rises} above $\theta_{00}$ as $y_0$ moves off the equator, by
$\Delta_y \sin^2(\pi y_0/2y_1)$: the SB08 family is not a rigid translation, and
a solstitial SB08 forcing is stronger as well as displaced. Seasonal cycles vary
$y_0$ in time and are available only for SB08. Under latent heating the profile
these formulae describe is $\theta_{\rm rad}$ and the contrast is $\dyr$.

\paragraph{Parameter values} Every run in the sweep uses
$g = 9.81$~m\,s$^{-2}$, $H = 16$~km, $d = 4$~km, $\beta = 2\times10^{-11}$
m$^{-1}$\,s$^{-1}$, $T_0 = 300$~K, $\theta_{00} = 330$~K, $y_1 = 9439$~km,
$\Delta_y = 50$~K, $\tau = 37$~d, $k_v = 7.786\times10^{5}$~m$^2$\,s$^{-1}$,
$\epsilon_u = 10^{-8}$~s$^{-1}$, $\kappa_\theta = 0$, and the production
formulation (staggered $v$, EMFD gate on, MC stencil, kinematic vertical-advection
gate, numba backend). The derived constants are
$C = 5.1619\times10^{6}$~J\,m$^{-2}$\,K$^{-1}$,
$\Lambda = 0.48431$~K per kg\,m$^{-2}$\,s$^{-1}$ and
$\Shat = 7.7429\times10^{7}$~J\,m$^{-2}$. Defaults for what the sweep varies:
$a = 0.85$, $W_c = 50$~kg\,m$^{-2}$, $\tau_c = 14400$~s,
$E_0 = 4.6\times10^{-5}$~kg\,m$^{-2}$\,s$^{-1}$ ($3.97$~\mmday),
$D = 10^{6}$~m$^2$\,s$^{-1}$, $v_d = 2.5$~\ms, $\Delta_z = 60$~K,
$\dyr = 75$~K.

\section{Run inventory}
\label{sec:inventory}

The sweep is defined by \path|scripts/v2_sweep_manifest.py|, which is the single
source of truth: it computes each run's parameters, records the timestep each
configuration needs, and generates the command line. The table below is
generated from it.

{\footnotesize
\begin{longtable}{llrrp{7.0cm}}
\caption{Every run in the sweep.}\label{tab:inventory}\\
\toprule
run & block & $n_y$ & $\Delta t$ (s) & what it varies\\
\midrule
\endfirsthead
\toprule
run & block & $n_y$ & $\Delta t$ (s) & what it varies\\
\midrule
\endhead
\bottomrule
\endfoot
\texttt{A\_a075\_wc30} & A & 801 & 30 & GMS plane: a=0.75, W\_c=30.0, r=0.495\\
\texttt{A\_a075\_wc40} & A & 801 & 30 & GMS plane: a=0.75, W\_c=40.0, r=0.656\\
\texttt{A\_a075\_wc50} & A & 801 & 30 & GMS plane: a=0.75, W\_c=50.0, r=0.818\\
\texttt{A\_a075\_wc60} & A & 801 & 30 & GMS plane: a=0.75, W\_c=60.0, r=0.979\\
\texttt{A\_a079\_wc30} & A & 801 & 30 & GMS plane: a=0.79, W\_c=30.0, r=0.574\\
\texttt{A\_a079\_wc40} & A & 801 & 30 & GMS plane: a=0.79, W\_c=40.0, r=0.761\\
\texttt{A\_a079\_wc50} & A & 801 & 30 & GMS plane: a=0.79, W\_c=50.0, r=0.949\\
\texttt{A\_a079\_wc60} & A & 801 & 30 & GMS plane: a=0.79, W\_c=60.0, r=1.136\\
\texttt{A\_a083\_wc30} & A & 801 & 30 & GMS plane: a=0.83, W\_c=30.0, r=0.653\\
\texttt{A\_a083\_wc40} & A & 801 & 30 & GMS plane: a=0.83, W\_c=40.0, r=0.867\\
\texttt{A\_a083\_wc50} & A & 801 & 30 & GMS plane: a=0.83, W\_c=50.0, r=1.080\\
\texttt{A\_a083\_wc60} & A & 801 & 15 & GMS plane: a=0.83, W\_c=60.0, r=1.293\\
\texttt{A\_a085\_wc30} & A & 801 & 30 & GMS plane: a=0.85, W\_c=30.0, r=0.693\\
\texttt{A\_a085\_wc40} & A & 801 & 30 & GMS plane: a=0.85, W\_c=40.0, r=0.919\\
\texttt{A\_a085\_wc50} & A & 801 & 30 & GMS plane: a=0.85, W\_c=50.0, r=1.145\\
\texttt{A\_a085\_wc60} & A & 801 & 15 & GMS plane: a=0.85, W\_c=60.0, r=1.371\\
\texttt{A\_a087\_wc30} & A & 801 & 30 & GMS plane: a=0.87, W\_c=30.0, r=0.733\\
\texttt{A\_a087\_wc40} & A & 801 & 30 & GMS plane: a=0.87, W\_c=40.0, r=0.972\\
\texttt{A\_a087\_wc50} & A & 801 & 30 & GMS plane: a=0.87, W\_c=50.0, r=1.210\\
\texttt{A\_a087\_wc60} & A & 801 & 15 & GMS plane: a=0.87, W\_c=60.0, r=1.449\\
\texttt{B\_a085\_wc35} & B & 801 & 30 & fine W\_c transect: W\_c=35.0, r=0.806\\
\texttt{B\_a085\_wc42} & B & 801 & 30 & fine W\_c transect: W\_c=42.0, r=0.964\\
\texttt{B\_a085\_wc44} & B & 801 & 30 & fine W\_c transect: W\_c=44.0, r=1.009\\
\texttt{B\_a085\_wc46} & B & 801 & 30 & fine W\_c transect: W\_c=46.0, r=1.055\\
\texttt{B\_a085\_wc55} & B & 801 & 30 & fine W\_c transect: W\_c=55.0, r=1.258\\
\texttt{C\_r085\_a075} & C & 801 & 30 & matched r=0.85: a=0.75 needs W\_c=51.99\\
\texttt{C\_r085\_a079} & C & 801 & 30 & matched r=0.85: a=0.79 needs W\_c=44.73\\
\texttt{C\_r085\_a083} & C & 801 & 30 & matched r=0.85: a=0.83 needs W\_c=39.23\\
\texttt{C\_r085\_a087} & C & 801 & 30 & matched r=0.85: a=0.87 needs W\_c=34.91\\
\texttt{C\_r105\_a075} & C & 801 & 30 & matched r=1.05: a=0.75 needs W\_c=64.38\\
\texttt{C\_r105\_a079} & C & 801 & 30 & matched r=1.05: a=0.79 needs W\_c=55.41\\
\texttt{C\_r105\_a083} & C & 801 & 30 & matched r=1.05: a=0.83 needs W\_c=48.61\\
\texttt{C\_r105\_a087} & C & 801 & 30 & matched r=1.05: a=0.87 needs W\_c=43.28\\
\texttt{D\_P\_dyr050} & D & 801 & 30 & radiative contrast 50.0 K on the P reference\\
\texttt{D\_P\_dyr060} & D & 801 & 30 & radiative contrast 60.0 K on the P reference\\
\texttt{D\_P\_dyr090} & D & 801 & 30 & radiative contrast 90.0 K on the P reference\\
\texttt{D\_P\_dyr100} & D & 801 & 30 & radiative contrast 100.0 K on the P reference\\
\texttt{D\_P\_dyr110} & D & 801 & 30 & radiative contrast 110.0 K on the P reference\\
\texttt{D\_N\_dyr050} & D & 801 & 30 & radiative contrast 50.0 K on the N reference\\
\texttt{D\_N\_dyr060} & D & 801 & 30 & radiative contrast 60.0 K on the N reference\\
\texttt{D\_N\_dyr090} & D & 801 & 30 & radiative contrast 90.0 K on the N reference\\
\texttt{D\_N\_dyr100} & D & 801 & 15 & radiative contrast 100.0 K on the N reference\\
\texttt{D\_N\_dyr110} & D & 801 & 15 & radiative contrast 110.0 K on the N reference\\
\texttt{E\_a085\_wc40\_tc001800} & E & 801 & 30 & tau\_c=1800 s at a=0.85, W\_c=40, r=0.906\\
\texttt{E\_a085\_wc40\_tc007200} & E & 801 & 30 & tau\_c=7200 s at a=0.85, W\_c=40, r=0.912\\
\texttt{E\_a085\_wc40\_tc043200} & E & 801 & 30 & tau\_c=43200 s at a=0.85, W\_c=40, r=0.949\\
\texttt{E\_a085\_wc40\_tc086400} & E & 801 & 30 & tau\_c=86400 s at a=0.85, W\_c=40, r=0.994\\
\texttt{E\_a085\_wc40\_tc172800} & E & 801 & 30 & tau\_c=172800 s at a=0.85, W\_c=40, r=1.084\\
\texttt{E\_P\_tc001800} & E & 801 & 30 & tau\_c=1800 s on the P reference\\
\texttt{E\_P\_tc172800} & E & 801 & 30 & tau\_c=172800 s on the P reference\\
\texttt{E\_N\_tc001800} & E & 801 & 30 & tau\_c=1800 s on the N reference\\
\texttt{E\_N\_tc172800} & E & 801 & 30 & tau\_c=172800 s on the N reference\\
\texttt{F\_P\_D0p00e00} & F & 801 & 30 & D=0 m\textasciicircum{}2/s on the P reference\\
\texttt{F\_P\_D2p50e05} & F & 801 & 30 & D=2.5e+05 m\textasciicircum{}2/s on the P reference\\
\texttt{F\_P\_D5p00e05} & F & 801 & 30 & D=5e+05 m\textasciicircum{}2/s on the P reference\\
\texttt{F\_P\_D2p00e06} & F & 801 & 30 & D=2e+06 m\textasciicircum{}2/s on the P reference\\
\texttt{F\_P\_D4p00e06} & F & 801 & 30 & D=4e+06 m\textasciicircum{}2/s on the P reference\\
\texttt{F\_N\_D0p00e00} & F & 801 & 30 & D=0 m\textasciicircum{}2/s on the N reference\\
\texttt{F\_N\_D2p50e05} & F & 801 & 30 & D=2.5e+05 m\textasciicircum{}2/s on the N reference\\
\texttt{F\_N\_D5p00e05} & F & 801 & 30 & D=5e+05 m\textasciicircum{}2/s on the N reference\\
\texttt{F\_N\_D2p00e06} & F & 801 & 30 & D=2e+06 m\textasciicircum{}2/s on the N reference\\
\texttt{F\_N\_D4p00e06} & F & 801 & 30 & D=4e+06 m\textasciicircum{}2/s on the N reference\\
\texttt{G\_P\_E2p30em05} & G & 801 & 30 & E\_0=2.3e-05 kg/m2/s (1.99 mm/day) on P\\
\texttt{G\_P\_E3p45em05} & G & 801 & 30 & E\_0=3.45e-05 kg/m2/s (2.98 mm/day) on P\\
\texttt{G\_P\_E6p90em05} & G & 801 & 30 & E\_0=6.9e-05 kg/m2/s (5.96 mm/day) on P\\
\texttt{G\_P\_E9p20em05} & G & 801 & 30 & E\_0=9.2e-05 kg/m2/s (7.95 mm/day) on P\\
\texttt{G\_N\_E2p30em05} & G & 801 & 30 & E\_0=2.3e-05 kg/m2/s (1.99 mm/day) on N\\
\texttt{G\_N\_E3p45em05} & G & 801 & 30 & E\_0=3.45e-05 kg/m2/s (2.98 mm/day) on N\\
\texttt{G\_N\_E6p90em05} & G & 801 & 30 & E\_0=6.9e-05 kg/m2/s (5.96 mm/day) on N\\
\texttt{G\_N\_E9p20em05} & G & 801 & 30 & E\_0=9.2e-05 kg/m2/s (7.95 mm/day) on N\\
\texttt{H\_P\_vd0p00} & H & 801 & 30 & v\_d=0.0 m/s on the P reference\\
\texttt{H\_P\_vd1p25} & H & 801 & 30 & v\_d=1.25 m/s on the P reference\\
\texttt{H\_P\_vd5p00} & H & 801 & 30 & v\_d=5.0 m/s on the P reference\\
\texttt{H\_P\_vd7p50} & H & 801 & 30 & v\_d=7.5 m/s on the P reference\\
\texttt{H\_N\_vd0p00} & H & 801 & 30 & v\_d=0.0 m/s on the N reference\\
\texttt{H\_N\_vd1p25} & H & 801 & 30 & v\_d=1.25 m/s on the N reference\\
\texttt{H\_N\_vd5p00} & H & 801 & 30 & v\_d=5.0 m/s on the N reference\\
\texttt{H\_N\_vd7p50} & H & 801 & 30 & v\_d=7.5 m/s on the N reference\\
\texttt{I\_P\_dz30} & I & 801 & 6 & Delta\_z=30.0 K on P, r=1.523\\
\texttt{I\_P\_dz45} & I & 801 & 30 & Delta\_z=45.0 K on P, r=1.015\\
\texttt{I\_P\_dz75} & I & 801 & 30 & Delta\_z=75.0 K on P, r=0.609\\
\texttt{I\_P\_dz90} & I & 801 & 30 & Delta\_z=90.0 K on P, r=0.508\\
\texttt{I\_N\_dz30} & I & 801 & 6 & Delta\_z=30.0 K on N, r=2.290\\
\texttt{I\_N\_dz45} & I & 801 & 15 & Delta\_z=45.0 K on N, r=1.527\\
\texttt{I\_N\_dz75} & I & 801 & 30 & Delta\_z=75.0 K on N, r=0.916\\
\texttt{I\_N\_dz90} & I & 801 & 30 & Delta\_z=90.0 K on N, r=0.763\\
\texttt{J\_P\_lam000} & J & 801 & 30 & Lambda = 0.0 x L\_v/C on the P reference\\
\texttt{J\_P\_lam025} & J & 801 & 30 & Lambda = 0.25 x L\_v/C on the P reference\\
\texttt{J\_P\_lam050} & J & 801 & 30 & Lambda = 0.5 x L\_v/C on the P reference\\
\texttt{J\_P\_lam075} & J & 801 & 30 & Lambda = 0.75 x L\_v/C on the P reference\\
\texttt{J\_P\_lam125} & J & 801 & 30 & Lambda = 1.25 x L\_v/C on the P reference\\
\texttt{J\_N\_lam000} & J & 801 & 30 & Lambda = 0.0 x L\_v/C on the N reference\\
\texttt{J\_N\_lam025} & J & 801 & 30 & Lambda = 0.25 x L\_v/C on the N reference\\
\texttt{J\_N\_lam050} & J & 801 & 30 & Lambda = 0.5 x L\_v/C on the N reference\\
\texttt{J\_N\_lam075} & J & 801 & 30 & Lambda = 0.75 x L\_v/C on the N reference\\
\texttt{J\_N\_lam125} & J & 801 & 15 & Lambda = 1.25 x L\_v/C on the N reference\\
\texttt{K\_wc40\_wi55} & K & 801 & 30 & a=0.85, W\_c=40.0, W(t=0)=55.0 (r=0.919)\\
\texttt{K\_wc40\_wi70} & K & 801 & 30 & a=0.85, W\_c=40.0, W(t=0)=70.0 (r=0.919)\\
\texttt{K\_wc42\_wi57} & K & 801 & 30 & a=0.85, W\_c=42.0, W(t=0)=57.0 (r=0.964)\\
\texttt{K\_wc42\_wi72} & K & 801 & 30 & a=0.85, W\_c=42.0, W(t=0)=72.0 (r=0.964)\\
\texttt{K\_wc44\_wi59} & K & 801 & 30 & a=0.85, W\_c=44.0, W(t=0)=59.0 (r=1.009)\\
\texttt{K\_wc44\_wi74} & K & 801 & 30 & a=0.85, W\_c=44.0, W(t=0)=74.0 (r=1.009)\\
\texttt{K\_wc46\_wi61} & K & 801 & 30 & a=0.85, W\_c=46.0, W(t=0)=61.0 (r=1.055)\\
\texttt{K\_wc46\_wi76} & K & 801 & 30 & a=0.85, W\_c=46.0, W(t=0)=76.0 (r=1.055)\\
\texttt{L\_P\_sin2\_y0700} & L & 801 & 30 & sin\textasciicircum{}2 forcing peak at y\_0=700 km on P\\
\texttt{L\_P\_sin2\_y1400} & L & 801 & 15 & sin\textasciicircum{}2 forcing peak at y\_0=1400 km on P\\
\texttt{L\_P\_sin2\_y2100} & L & 801 & 15 & sin\textasciicircum{}2 forcing peak at y\_0=2100 km on P\\
\texttt{L\_P\_sb08\_y0000} & L & 801 & 30 & SB08 forcing peak at y\_0=0 km on P\\
\texttt{L\_P\_sb08\_y0700} & L & 801 & 30 & SB08 forcing peak at y\_0=700 km on P\\
\texttt{L\_P\_sb08\_y1400} & L & 801 & 15 & SB08 forcing peak at y\_0=1400 km on P\\
\texttt{L\_P\_sb08\_y2100} & L & 801 & 15 & SB08 forcing peak at y\_0=2100 km on P\\
\texttt{L\_P\_sb08\_y2800} & L & 801 & 6 & SB08 forcing peak at y\_0=2800 km on P\\
\texttt{L\_N\_sin2\_y0700} & L & 801 & 15 & sin\textasciicircum{}2 forcing peak at y\_0=700 km on N\\
\texttt{L\_N\_sin2\_y1400} & L & 801 & 15 & sin\textasciicircum{}2 forcing peak at y\_0=1400 km on N\\
\texttt{L\_N\_sin2\_y2100} & L & 801 & 15 & sin\textasciicircum{}2 forcing peak at y\_0=2100 km on N\\
\texttt{L\_N\_sb08\_y0000} & L & 801 & 30 & SB08 forcing peak at y\_0=0 km on N\\
\texttt{L\_N\_sb08\_y0700} & L & 801 & 15 & SB08 forcing peak at y\_0=700 km on N\\
\texttt{L\_N\_sb08\_y1400} & L & 801 & 15 & SB08 forcing peak at y\_0=1400 km on N\\
\texttt{L\_N\_sb08\_y2100} & L & 801 & 6 & SB08 forcing peak at y\_0=2100 km on N\\
\texttt{L\_N\_sb08\_y2800} & L & 801 & 6 & SB08 forcing peak at y\_0=2800 km on N\\
\texttt{M\_P\_amp0350\_p360} & M & 801 & 15 & seasonal amplitude 350 km, 360-day period, P\\
\texttt{M\_P\_amp0700\_p360} & M & 801 & 15 & seasonal amplitude 700 km, 360-day period, P\\
\texttt{M\_P\_amp1400\_p360} & M & 801 & 15 & seasonal amplitude 1400 km, 360-day period, P\\
\texttt{M\_P\_amp2100\_p360} & M & 801 & 15 & seasonal amplitude 2100 km, 360-day period, P\\
\texttt{M\_P\_amp2800\_p360} & M & 801 & 15 & seasonal amplitude 2800 km, 360-day period, P\\
\texttt{M\_N\_amp0350\_p360} & M & 801 & 15 & seasonal amplitude 350 km, 360-day period, N\\
\texttt{M\_N\_amp0700\_p360} & M & 801 & 15 & seasonal amplitude 700 km, 360-day period, N\\
\texttt{M\_N\_amp1400\_p360} & M & 801 & 15 & seasonal amplitude 1400 km, 360-day period, N\\
\texttt{M\_N\_amp2100\_p360} & M & 801 & 15 & seasonal amplitude 2100 km, 360-day period, N\\
\texttt{M\_N\_amp2800\_p360} & M & 801 & 15 & seasonal amplitude 2800 km, 360-day period, N\\
\texttt{M\_P\_amp1400\_p0090} & M & 801 & 15 & seasonal period 90 days at amplitude 1400 km, P\\
\texttt{M\_P\_amp1400\_p0180} & M & 801 & 15 & seasonal period 180 days at amplitude 1400 km, P\\
\texttt{M\_P\_amp1400\_p0720} & M & 801 & 15 & seasonal period 720 days at amplitude 1400 km, P\\
\texttt{M\_P\_amp1400\_p1440} & M & 801 & 15 & seasonal period 1440 days at amplitude 1400 km, P\\
\texttt{M\_N\_amp1400\_p0090} & M & 801 & 15 & seasonal period 90 days at amplitude 1400 km, N\\
\texttt{M\_N\_amp1400\_p1440} & M & 801 & 15 & seasonal period 1440 days at amplitude 1400 km, N\\
\texttt{M\_P\_amp1400\_p360\_tanh} & M & 801 & 15 & tanh seasonal waveform, amplitude 1400 km, 360-day period, P\\
\texttt{M\_P\_amp1400\_p360\_square} & M & 801 & 15 & square seasonal waveform, amplitude 1400 km, 360-day period, P\\
\texttt{N\_P\_v1} & N & 801 & 30 & moist V1 twin of the P reference: no latent heating, no retarget, thermodynamic gate\\
\texttt{N\_N\_v1} & N & 801 & 30 & moist V1 twin of the N reference: no latent heating, no retarget, thermodynamic gate\\
\texttt{O\_P\_ny1601} & O & 1601 & 15 & P reference at ny=1601, dt=15\\
\texttt{O\_P\_ny401} & O & 401 & 60 & P reference at ny=401, dt=60\\
\texttt{O\_P\_dt15} & O & 801 & 15 & P reference at ny=801, dt=15\\
\texttt{O\_N\_ny1601} & O & 1601 & 15 & N reference at ny=1601, dt=15\\
\texttt{O\_N\_ny401} & O & 401 & 60 & N reference at ny=401, dt=60\\
\texttt{O\_N\_dt15} & O & 801 & 15 & N reference at ny=801, dt=15\\
\texttt{O\_crit\_ny1601} & O & 1601 & 15 & near-critical a=0.85, W\_c=44 at ny=1601, dt=15\\
\texttt{O\_P\_sb08\_y0700\_ny1601} & O & 1601 & 15 & off-equatorial SB08 y\_0=700 km, P, at ny=1601, dt=15\\
\texttt{O\_P\_amp1400\_p360\_ny1601} & O & 1601 & 15 & seasonal amplitude 1400 km, P, at ny=1601, dt=15\\
\texttt{O\_P\_sb08\_y0700\_ny0801\_dt6} & O & 801 & 6 & off-equatorial SB08 y\_0=700 km, P, at ny=801, dt=6, 2500 d\\
\texttt{O\_P\_sb08\_y0700\_ny1601\_dt6} & O & 1601 & 6 & off-equatorial SB08 y\_0=700 km, P, at ny=1601, dt=6, 2500 d\\
\end{longtable}
}


\section{Commands}
\label{sec:repro}

From the repository root, with the \texttt{claude-swm} conda environment:

\begin{verbatim}
# what the sweep is, and what it costs
python scripts/v2_sweep_manifest.py

# stability screen: every configuration for 100 days, 16 at a time
python scripts/v2_sweep_run.py --smoke 100 --nproc 16

# the sweep itself (about 2.5 h at 20 concurrent on berimbau)
python scripts/v2_sweep_run.py --nproc 20

# re-run only what failed or is missing
python scripts/v2_sweep_run.py --nproc 20 --skip-done

# diagnostics, every figure, and every number quoted above
python scripts/v2_sweep_figures.py
python scripts/v2_sweep_report_tables.py

# this document
latexmk -pdf -outdir=docs docs/moist_v2_sweep_report.tex
\end{verbatim}

Output lives in \path|model_output/moist_v2_sweep/|, one directory per run
holding \path|out.nc| and a final restart file, with \path|logs/| alongside and
\path|status.jsonl| recording each run's exit state, wall time and stored days.
The scalar table for every run is \path|all_runs.csv|, the quoted statistics are
\path|numbers.json|, and the figures are in \path|figs/|.

\end{document}
